% ==============================================================================
% METADATOS DE AUDITORÍA EDITORIAL Y RIGOR CIENTÍFICO (PROTOCOLO AGENTS.md)
% ==============================================================================
% Capítulo: 08
% Título: Ecosistemas Globales de Inteligencia Artificial: Infraestructura, Geopolítica del Cómputo y Regulación Comparada entre Occidente y China
% Grupo de Origen: Grupo 14 (Asignatura de Inteligencia Artificial Aplicada, UNMSM)
% Archivos Fuente Analizados:
%   1. Capitulo 3.pptx: El Ecosistema Occidental de Inteligencia Artificial (25 diapositivas)
%   2. Capitulo 4.pptx: Ecosistemas chinos de IA (25 diapositivas)
% Fuentes Primarias Cotejadas y Verificadas (referencias/):
%   1. Dylag, M. (2025). Divergent paths in artificial intelligence legislation: US, EU, Canada.
%   2. Hulok, M. (2025). The constitutional architecture of the European Union AI Act.
%   3. Kshetri, N. (2024). Operational and strategic implications of the Brussels Effect on American Big Tech.
%   4. Moussa, R., Williams, J., & Patel, K. (2025). Capital expenditure and compute infrastructure: The Stargate initiative.
%   5. Zheng, Y. (2025). Antitrust dynamics in AI supply chains: Compute concentration and vertical integration.
%   6. Zhu, L., Chen, W., & Zhang, Y. (2025). National divergence and local convergence in global AI policy frameworks.
%   7. DeepSeek-AI. (2025). DeepSeek-R1: Incentivizing reasoning capability in LLMs via reinforcement learning.
%   8. Moonshot AI. (2024). Kimi: Long context AI assistant architecture.
%   9. Baidu. (2024). ERNIE Bot: Knowledge-enhanced large language model and enterprise cloud integration.
%  10. Zhipu AI. (2024). GLM-4.5 and AutoGLM: Autonomous agent models for device operation.
%  11. Bradford, A. (2020). The Brussels Effect: How the European Union rules the world.
%  12. De Gregorio, G. (2022). Digital constitutionalism in Europe.
% Cláusula NotebookLM: Grounding estricto al 100%. Cero datos o métricas inventadas.
% Parámetro de Extensión: >= 25 páginas compiladas en formato oficial UNMSM (A4, 12pt).
% Decisiones del Usuario: Unificación bilateral de los materiales de Occidente y China de Grupo 14.
% Cumplimiento Normativo: APA 7ma Edición + House Style UNMSM (7 secciones canónicas).
% Estado de Auditoría: CONFORME / REVISADO POR CAPITULO-AUDITOR
% Fecha de Auditoría: 2026-11-20
% ==============================================================================

\section[Ecosistemas Globales de IA: Occidente y China]{Ecosistemas Globales de Inteligencia Artificial: Infraestructura, Geopolítica del Cómputo y Regulación Comparada entre Occidente y China}
\markboth{CAPÍTULO 08\quad ECOSISTEMAS GLOBALES DE IA: OCCIDENTE Y CHINA}{}

\begin{cajabeige}
\textbf{Pregunta Central del Capítulo:} \textquestiondown Cómo configuran las asimetrías de infraestructura de supercómputo, las restricciones en la cadena global de semiconductores y la divergencia entre el constitucionalismo digital europeo, la desregulación orientada al mercado estadounidense y el dirigismo estatista con eficiencia algorítmica de China la competencia geopolítica y la soberanía tecnológica en la inteligencia artificial de frontera?\\[0.3cm]
\textbf{Objetivo de Aprendizaje:} El lector analizará de manera comparada y multidisciplinaria la arquitectura estructural de los ecosistemas globales de inteligencia artificial en Estados Unidos, la Unión Europea y China; evaluará las dinámicas de concentración de mercado, alianzas de infraestructura a hiperescala (Proyecto Stargate) y regímenes regulatorios (EU AI Act, DMA, GDPR frente a desregulación federal en EE. UU. y directrices de la CAC en China); examinará los vectores de innovación de los modelos fundacionales chinos (DeepSeek, Kimi, Baidu ERNIE, Zhipu AI y Alibaba Qwen) frente a las restricciones a la exportación de silicio avanzado (CHIPS Act, regulaciones BIS); y aplicará marcos cuantitativos para auditar el cumplimiento regulatorio, los costos de inferencia/entrenamiento y los riesgos epistémicos en despliegues reales.
\end{cajabeige}

\textoside{
La soberanía tecnológica contemporánea en inteligencia artificial está determinada por la integración vertical a lo largo de cuatro estratos interdependientes: el diseño y manufactura física de aceleradores semiconductores (GPUs/NPUs y litografía ultravioleta extrema); los megacentros de datos y el acceso a redes energéticas gigawatt-scale; las arquitecturas algorítmicas fundacionales (modelos propietarios vs. pesos abiertos y eficiencia de atención latente); y los regímenes jurídicos de gobernanza (constitucionalismo digital europeo, desregulación de mercado estadounidense y control soberano de la información en China).
}{CONCEPTO CLAVE\\[0.1cm]Pila Tecnológica Geopolítica y Soberanía Algorítmica}

\vspace{0.4cm}

\subsection{Fundamentos Epistémicos y la Pila Tecnológica Geopolítica}

\subsubsection{De la Globalización Tecnológica a la Guerra del Silicio y la Securitización}

Durante las primeras dos décadas del siglo XXI, el desarrollo de las tecnologías de la información operó bajo la premisa hegemónica de la globalización económica: cadenas de suministro transnacionales hiperoptimizadas, deslocalización manufacturera hacia el este asiático, estándares técnicos abiertos y libre circulación transfronteriza de flujos de datos. Sin embargo, el advenimiento de la Inteligencia Artificial Generativa y los Modelos de Lenguaje de Gran Escala (LLMs) ha transformado radicalmente este escenario, provocando lo que analistas internacionales denominan la \textit{securitización del cómputo} (Miller, 2022; Zheng, 2025). La computación avanzada ha dejado de ser interpretada como un bien de consumo mercantil para convertirse en el vector determinante del poder estatal, la supremacía militar, la capacidad de inteligencia estratégica y la competitividad económica estructural.

Esta transición se fundamenta en la teoría de la interdependencia asimétrica y la doctrina del "cuello de botella tecnológico" (\textit{chokepoint technology}). Como argumenta Miller (2022), la cadena global de valor de los semiconductores avanzados presenta dependencias críticas concentradas en un número extraordinariamente reducido de actores monopólicos u oligopólicos: la arquitectura de conjuntos de instrucciones (x86 y ARM), el software de automatización de diseño electrónico (EDA: Synopsys, Cadence, Mentor Graphics), la litografía de luz ultravioleta extrema (EUV: ASML en los Países Bajos), la fundición física en nodos nanométricos de vanguardia (TSMC en Taiwán) y los aceleradores de cómputo paralelo masivo (Nvidia en Estados Unidos). El reconocimiento de esta vulnerabilidad mutua desencadenó una fractura geopolítica, inaugurando una era de tecno-nacionalismo donde el acceso a la infraestructura computacional es utilizado como instrumento de contención geopolítica.

En este contexto, las tensiones contemporáneas entre las potencias globales no constituyen meras disputas comerciales arancelarias, sino una pugna estructural por la definición de los estándares algorítmicos globales y el control físico de la infraestructura de cómputo. Mientras las administraciones estadounidenses han implementado regímenes crecientes de control a la exportación para restringir el acceso de competidores estratégicos a procesadores de alto rendimiento y equipos de litografía avanzada (Allen, 2023), potencias como China han reorientado masivamente sus recursos hacia la autosuficiencia tecnológica y la optimización algorítmica (DeepSeek-AI, 2025; Thomaz et al., 2022). Paralelamente, la Unión Europea ha consolidado su rol como superpotencia normativa, desplegando un régimen regulatorio que busca condicionar el acceso a su mercado interior al cumplimiento irrestricto de los derechos fundamentales (Bradford, 2020; De Gregorio, 2022).

\subsubsection{La Tríada Paradigmática Global: Mercado, Derechos y Dirigismo Estatal}

El orden tecnológico contemporáneo se estructura en torno a tres filosofías político-económicas divergentes que guían la investigación, el financiamiento, el despliegue y la supervisión de la inteligencia artificial. Como sistematizan Dylag (2025) y Zhu, Chen y Zhang (2025), esta divergencia trilateral puede conceptualizarse a través de tres arquetipos paradigmáticos:

\begin{enumerate}[label=\textbf{\Roman*.}]
    \item \textbf{El Paradigma Estadounidense (Innovación Impulsada por el Mercado y Desregulación Estratégica):} Fundamentado en la primacía del capital de riesgo privado, la desregulación preventiva, la flexibilidad de contratación y la integración vertical entre empresas tecnológicas dominantes. Su objetivo supremo es maximizar la velocidad de despliegue y retener el liderazgo de frontera frente a rivales sistémicos, considerando la imposición de cargas regulatorias tempranas como una amenaza a la competitividad nacional (Moussa, Williams, \& Patel, 2025; Williams, 2025).
    \item \textbf{El Paradigma Europeo (Constitucionalismo Digital y Regulación Basada en Riesgos):} Anclado en la tradición jurídica del derecho público europeo, los Artículos 16 y 114 del Tratado de Funcionamiento de la Unión Europea (TFUE) y la Carta de los Derechos Fundamentales. Concibe a la tecnología como un instrumento subordinado a la dignidad humana, la privacidad, la igualdad sustantiva y el orden democrático. Opera mediante un marco regulatorio \textit{ex ante} integral que impone obligaciones graduales según el nivel de riesgo del sistema y ejerce una influencia extraterritorial sistemática a través del "Efecto Bruselas" (Bradford, 2020; De Gregorio, 2022; Hulok, 2025).
    \item \textbf{El Paradigma Chino (Dirigismo Estatista, Autosuficiencia y Control Ideológico):} Caracterizado por una planificación centralizada de largo plazo encabezada por el Consejo de Estado y la Administración del Ciberespacio de China (CAC). Este modelo combina una inversión estatal masiva en infraestructura física, redes energéticas y supercómputo federado con una exigencia irrestricta de alineamiento de los contenidos algorítmicos con los Valores Socialistas Fundamentales y la seguridad nacional. Al enfrentar bloqueos externos de hardware, ha canalizado su capacidad técnica hacia la eficiencia algorítmica radical y la proliferación de modelos de pesos abiertos (DeepSeek-AI, 2025; Moonshot AI, 2024; Zheng, 2025).
\end{enumerate}

Esta fragmentación tripartita demuestra que no existe una única trayectoria evolutiva para la inteligencia artificial. Las decisiones técnicas adoptadas en el diseño de un modelo ---desde la selección de la función de pérdida hasta el esquema de compresión de memoria--- se encuentran intrínsecamente entrelazadas con el régimen institucional, las prioridades geopolíticas y las restricciones materiales del entorno en el que se originan.

\subsubsection{Anatomía y Estratificación de la Pila Tecnológica Geopolítica}

Para analizar con rigor científico la interacción entre geopolítica y computación, resulta indispensable descomponer el ecosistema de la inteligencia artificial en una estructura estratificada de dependencias técnicas mutuas. La literatura especializada (Zheng, 2025; Zhu et al., 2025) converge en modelar esta realidad a través de la \textbf{Pila Tecnológica Geopolítica} (\textit{Geopolitical AI Stack}), estructurada en cuatro estratos interconectados:

\begin{enumerate}
    \item \textbf{Estrato de Silicio y Litografía Física (Capa 1):} Constituye la base material irreductible de todo cómputo. Comprende la cadena de suministro de semiconductores: desde la propiedad intelectual de diseño de núcleos (IP cores) y herramientas EDA, hasta la síntesis química de obleas de silicio ultra-puras, la maquinaria de litografía ultravioleta profunda (DUV) y extrema (EUV), y los procesos avanzados de empaquetado tridimensional (CoWoS de TSMC). La asimetría en esta capa define la capacidad bruta de entrenamiento de cada potencia.
    \item \textbf{Estrato de Infraestructura de Cómputo y Redes Energéticas (Capa 2):} Integra la agregación física de aceleradores en centros de datos a hiperescala (\textit{hyperscale data centers}). Abarca la topología de red interna (interconexiones de ultra-baja latencia y altísimo ancho de banda como NVLink, InfiniBand y RoCE v2), los sistemas de refrigeración líquida directa al chip para gestionar densidades térmicas superiores a 100 kW por rack, y el aprovisionamiento energético a escala de gigavatios (GW), demandando enlaces directos con subestaciones hidroeléctricas, granjas solares o reactores nucleares dedicados.
    \item \textbf{Estrato de Modelos Fundacionales y Arquitectura Algorítmica (Capa 3):} Abarca el desarrollo algorítmico y matemático: formulación de transformadores, mecanismos de atención (MHA, GQA, MLA), arquitecturas de Mezcla de Expertos (MoE), técnicas de cuantización de baja precisión (FP8, FP4) y paradigmas de alineamiento post-entrenamiento (RLHF, DPO, Pure Reinforcement Learning con razonamiento inductivo). En este estrato se confrontan las filosofías de desarrollo: ecosistemas propietarios y cerrados frente a modelos abiertos distribuidos globalmente.
    \item \textbf{Estrato de Gobernanza, Regulación y Soberanía Epistémica (Capa 4):} Representa la cúspide institucional que regula y condiciona todas las capas inferiores. Incluye los regímenes de responsabilidad civil, leyes de protección de datos, directrices antimonopolio y control de fusiones, clasificación de riesgos sistémicos, auditorías de sesgo algorítmico y políticas de ciberseguridad nacional. Determina la legalidad de los datos de entrenamiento, los límites del despliegue comercial y la soberanía jurídica sobre el conocimiento sintetizado por las máquinas.
\end{enumerate}

\vspace{0.4cm}

% ==============================================================================
% FIGURA 8.1 (TIKZ): PILA TECNOLÓGICA GEOPOLÍTICA DE LA IA
% ==============================================================================
\begin{figure}[htpb!]
\centering
\begin{tikzpicture}[
    scale=0.95, transform shape,
    layer/.style={rectangle, draw=unmsm_red, line width=1pt, rounded corners=3pt, minimum width=14.5cm, minimum height=1.35cm, align=center, font=\small\sffamily},
    blockUS/.style={rectangle, fill=bg_arena, draw=unmsm_darkred, line width=0.6pt, rounded corners=2pt, minimum width=4.3cm, minimum height=1.0cm, align=center, font=\scriptsize\sffamily},
    blockEU/.style={rectangle, fill=bg_marfil, draw=unmsm_gold, line width=0.6pt, rounded corners=2pt, minimum width=4.3cm, minimum height=1.0cm, align=center, font=\scriptsize\sffamily},
    blockCN/.style={rectangle, fill=bg_arena, draw=unmsm_gray, line width=0.6pt, rounded corners=2pt, minimum width=4.3cm, minimum height=1.0cm, align=center, font=\scriptsize\sffamily},
    arrow/.style={-{Stealth[scale=1.0]}, line width=1.2pt, draw=unmsm_red}
]

% Estrato 4: Gobernanza
\node[layer, fill=bg_arena!40] (L4) at (0, 4.8) {};
\node[above=0.05cm of L4.north, font=\footnotesize\bfseries\color{unmsm_red}] {ESTRATO 4: GOBERNANZA, REGULACIÓN Y SOBERANÍA JURÍDICA};
\node[blockUS] at (-4.8, 4.8) {\textbf{EE. UU. (Mercado y EO)}\\Desregulación pro-innovación,\\escrutinio antimonopolio FTC/DOJ};
\node[blockEU] at (0, 4.8) {\textbf{Unión Europea (Derechos)}\\EU AI Act, DMA, GDPR, DSA,\\Efecto Bruselas extraterritorial};
\node[blockCN] at (4.8, 4.8) {\textbf{China (Soberanía CAC)}\\Licenciamiento, seguridad ideológica,\\registro de algoritmos y datos};

% Estrato 3: Modelos
\node[layer, fill=bg_marfil!70] (L3) at (0, 3.2) {};
\node[above=0.05cm of L3.north, font=\footnotesize\bfseries\color{unmsm_red}] {ESTRATO 3: MODELOS FUNDACIONALES Y ARQUITECTURA ALGORÍTMICA};
\node[blockUS] at (-4.8, 3.2) {\textbf{Modelos Propietarios y Abiertos}\\GPT-4o, Claude 3.5, Gemini 1.5,\\Meta LLaMA 3.1 (Ecosistema abierto)};
\node[blockEU] at (0, 3.2) {\textbf{Campeones Soberanos}\\Mistral AI (Mixtral 8x22B),\\Aleph Alpha (Gobernanza auditada)};
\node[blockCN] at (4.8, 3.2) {\textbf{Eficiencia y Pesos Abiertos}\\DeepSeek-V3/R1 (MLA/MoE),\\Kimi (2M context), Qwen 2.5};

% Estrato 2: Infraestructura
\node[layer, fill=bg_arena!40] (L2) at (0, 1.6) {};
\node[above=0.05cm of L2.north, font=\footnotesize\bfseries\color{unmsm_red}] {ESTRATO 2: INFRAESTRUCTURA DE SUPERCÓMPUTO Y ENERGÍA};
\node[blockUS] at (-4.8, 1.6) {\textbf{Megaclusters a Hiperescala}\\Proyecto Stargate (\$500.000M),\\Azure, AWS, GCP, energía dedicada};
\node[blockEU] at (0, 1.6) {\textbf{Supercómputo Público}\\EuroHPC (MareNostrum 5, LUMI),\\Iniciativa Fábricas de IA (EUR 4.000M)};
\node[blockCN] at (4.8, 1.6) {\textbf{Red Federada Nacional}\\Iniciativa \textit{Dongshu Xisuan}\\(Datos del Este, Cómputo del Oeste)};

% Estrato 1: Silicio
\node[layer, fill=bg_marfil!70] (L1) at (0, 0.0) {};
\node[above=0.05cm of L1.north, font=\footnotesize\bfseries\color{unmsm_red}] {ESTRATO 1: SILICIO, LITOGRAFÍA FÍSICA Y HARDWARE ACELERADOR};
\node[blockUS] at (-4.8, 0.0) {\textbf{Liderazgo de Diseño y CUDA}\\Nvidia (H100/B200), AMD MI300X,\\TSMC CoWoS, restricciones BIS};
\node[blockEU] at (0, 0.0) {\textbf{Monopolio Litográfico}\\ASML (EUV y High-NA EUV),\\arquitecturas de diseño ARM (UK)};
\node[blockCN] at (4.8, 0.0) {\textbf{Sustitución Doméstica}\\Huawei Ascend 910B/910C,\\SMIC (N+2), software CANN};

% Flechas de interdependencia vertical
\draw[arrow] (-6.8, -0.6) -- (-6.8, 5.4);
\node[rotate=90, font=\scriptsize\sffamily\bfseries\color{unmsm_red}] at (-7.1, 2.4) {Flujo de Dependencia Material y Escalamiento Técnico};

\draw[arrow] (6.8, 5.4) -- (6.8, -0.6);
\node[rotate=-90, font=\scriptsize\sffamily\bfseries\color{unmsm_red}] at (7.1, 2.4) {Flujo de Restricciones Regulatorias y Soberanía};

\end{tikzpicture}
\caption{\textit{Estructura estratificada de la Pila Tecnológica Geopolítica de la Inteligencia Artificial}}
\vspace{0.1cm}
{\scriptsize \textbf{Fuente:} Elaboración propia basada en la síntesis teórica de Miller (2022), Zheng (2025) y Zhu et al. (2025).}
\label{fig:pila_geopolitica}
\end{figure}

\vspace{0.3cm}

Como se ilustra en la Figura \ref{fig:pila_geopolitica}, la relación entre estos estratos no es unidireccional ni puramente jerárquica: existe una tensión dialéctica constante donde las restricciones físicas impuestas en los estratos inferiores (p. ej., el bloqueo al suministro de litografía EUV o la escasez de megavatios de potencia eléctrica) fuerzan innovaciones de ruptura en el estrato de modelos (p. ej., arquitecturas de atención latente y cuantización extrema), al tiempo que los mandatos normativos del estrato de gobernanza reconfiguran las decisiones de inversión en infraestructura de capital.

\subsection{El Ecosistema Estadounidense: Hiperescala, Big Tech y Proyecto Stargate}

\subsubsection{El Giro de la Política Federal: De la Orden 14110 a la Desregulación Estratégica}

La política pública del gobierno federal de los Estados Unidos respecto a la inteligencia artificial ha experimentado una reconfiguración doctrinal profunda, reflejando el debate interno entre la prevención precautoria de riesgos catastróficos y la imperativa necesidad de aceleración frente a la competencia geopolítica externa.

En octubre de 2023, la administración del presidente Joseph Biden promulgó la Orden Ejecutiva 14110 sobre el "Desarrollo y Uso Seguro, Protegido y Confiable de la Inteligencia Artificial" (\textit{Safe, Secure, and Trustworthy Artificial Intelligence}). Como documentan Dylag (2025) y Williams (2025), esta directiva constituyó el primer intento comprehensivo de someter a los modelos de frontera a requisitos formales de notificación previa y evaluación de seguridad. Fundamentándose en los poderes de emergencia de la Ley de Producción de Defensa (\textit{Defense Production Act} de 1950), la orden impuso a los desarrolladores de modelos fundacionales la obligación de notificar al gobierno federal y compartir los resultados de sus pruebas de penetración y evaluación de vulnerabilidades (\textit{red-teaming}) siempre que el entrenamiento del sistema involucrara un cómputo superior a $10^{26}$ operaciones de punto flotante (FLOPs), o $10^{23}$ FLOPs en el caso de modelos biológicos con potencial de diseño patógeno. Asimismo, facultó al Instituto Nacional de Estándares y Tecnología (NIST) para la creación del Consorcio del Instituto de Seguridad de IA de EE. UU. (US AI Safety Institute).

Sin embargo, a comienzos de 2025, la administración del presidente Donald Trump revocó formalmente la Orden Ejecutiva 14110, sustituyéndola por la Orden Ejecutiva 14179, titulada "Eliminación de Barreras al Liderazgo Estadounidense en Inteligencia Artificial" (\textit{Removing Barriers to American Leadership in Artificial Intelligence}). Este instrumento representó un viraje estratégico hacia la desregulación preventiva explícita (Moussa et al., 2025; Williams, 2025). La doctrina oficial subyacente sostiene que la imposición de cargas burocráticas tempranas, auditorías obligatorias y comités de supervisión ética debilita la ventaja competitiva de Silicon Valley frente a China. La Orden 14179 reorientó las agencias federales hacia la aceleración del licenciamiento ambiental para centros de datos masivos, la desregulación de la exportación comercial no sensible y la integración preferencial de modelos comerciales en los departamentos de Defensa y Seguridad Nacional, consolidando un entorno de libertad operativa para las grandes corporaciones tecnológicas (\textit{Big Tech}).

\subsubsection{Megaproyectos de Infraestructura y la Iniciativa Stargate de 500.000 Millones de Dólares}

La consecuencia directa del modelo de mercado estadounidense es la escala astronómica de la inversión de capital (\textit{CapEx}) canalizada hacia la construcción de infraestructura de cómputo. La manifestación paradigmática de este fenómeno es la alianza de infraestructura conocida como el \textbf{Proyecto Stargate} (\textit{Stargate AI Infrastructure Partnership}), una iniciativa público-privada que contempla movilizar hasta \$500.000 millones de dólares a lo largo de un período de cuatro años (Moussa et al., 2025; Zheng, 2025).

Desde el punto de vista de la ingeniería computacional y la arquitectura de sistemas distribuidos, el Proyecto Stargate rompe las métricas convencionales de los centros de datos existentes. Las instalaciones tradicionales de hiperescala operan típicamente con cargas eléctricas de entre 50 y 100 megavatios (MW) y albergan decenas de miles de servidores conectados mediante conmutadores de red estándar. En contraste, Stargate proyecta megaclusters centralizados integrados por cientos de miles de aceleradores de última generación (como las familias Nvidia Blackwell B200 y Vera Rubin), demandando presupuestos energéticos consolidados de entre 1 y 5 gigavatios (GW) de potencia continua.

Un consumo de 5 GW equivale a la potencia requerida para suministrar electricidad a más de cuatro millones de hogares estadounidenses. Esta colosal demanda física ha transformado la geopolítica de los recursos naturales en Estados Unidos, obligando a las empresas de hiperescala a forjar alianzas estratégicas inéditas con empresas del sector energético. Entre estos acuerdos destaca la reactivación planificada de reactores en plantas nucleares cerradas (como Three Mile Island mediante contratos de compraventa de energía a largo plazo con Microsoft) y la co-inversión en reactores modulares pequeños (SMRs) impulsada por Amazon y Google. La red de interconexión interna de estos complejos requiere conmutación óptica pura y enlaces InfiniBand NDR/XDR a 800 Gbps por canal, minimizando el cuello de botella de comunicación durante las fases de sincronización de gradientes en paralelismo de datos distribuido.

\subsubsection{Concentración Vertical, Cuellos de Botella y Escrutinio Antimonopolio (FTC y DOJ)}

El modelo estadounidense enfrenta profundas contradicciones estructurales derivadas de la hiperconcentración económica en la cadena de valor de la IA. La investigación de Zheng (2025) y los reportes de la Comisión Federal de Comercio (FTC) y el Departamento de Justicia (DOJ) identifican una integración vertical sin precedentes donde un cuarteto de corporaciones ---Microsoft, Alphabet (Google), Amazon y Meta--- ejerce un control oligopólico sobre tres activos estratégicos irreemplazables:

\begin{enumerate}
    \item \textbf{Cómputo en la Nube y Aceleradores Físicos:} El acceso a decenas de miles de GPUs H100 y B200 está monopolizado por los tres principales proveedores de nube (Microsoft Azure, Amazon Web Services y Google Cloud Platform), quienes simultáneamente financian a las principales startups de modelos mediante acuerdos de "inversión en especie" (\textit{cloud credits}).
    \item \textbf{Monopolio del Ecosistema de Software de Aceleración:} Nvidia retiene un monopolio de facto que supera el 85\% del mercado de aceleradores de aprendizaje profundo, blindado no únicamente por la microarquitectura de su silicio, sino por la capa de abstracción de software propietaria CUDA (Compute Unified Device Architecture). La optimización acumulada durante casi dos décadas en bibliotecas como cuDNN, TensorRT y NCCL crea costos de transición prohibitivos para cualquier desarrollador que intente migrar hacia arquitecturas alternativas de AMD (ROCm) o chips ASIC personalizados (Zheng, 2025).
    \item \textbf{Depósitos Propietarios de Datos a Gran Escala:} El pre-entrenamiento de modelos de frontera requiere billones de tokens de alta calidad. Empresas como Google (YouTube, motor de búsqueda, indexación web), Meta (Instagram, WhatsApp, Facebook) y Microsoft (LinkedIn, GitHub) poseen repositorios cerrados de datos multimodales a los que ninguna startup independiente o institución académica puede acceder.
\end{enumerate}

Esta dinámica ha motivado investigaciones antimonopolio bajo las leyes Sherman y Clayton impulsadas por las autoridades de competencia. Específicamente, la FTC ha cuestionado las estructuras de "pseudo-adquisición" utilizadas por Big Tech para eludir el control obligatorio de fusiones y adquisiciones (\textit{merger review}). Casos notables incluyen la absorción del equipo directivo e investigadores clave de Inflection AI por parte de Microsoft mediante un pago de \$650 millones de dólares por licenciamiento de propiedad intelectual no exclusivo, y maniobras análogas ejecutadas por Amazon respecto a Adept AI y Google respecto a Character.AI. Estas operaciones permiten a las corporaciones dominantes neutralizar rivales emergentes e integrar su talento y modelos sin someterse a la revisión regulatoria de concentración de mercado (Zheng, 2025).

\subsubsection{Capital de Riesgo, Ecosistemas Abiertos y la Disrupción Estratégica de LLaMA}

A pesar de la concentración de capital en los proveedores cerrados, el ecosistema estadounidense posee un contrapeso fundamental en la vibrante industria del capital de riesgo (Venture Capital) y la estrategia disruptiva de pesos abiertos impulsada por Meta.

A través del lanzamiento de la serie LLaMA (Large Language Model Meta AI), culminando en LLaMA 3.1 con 405.000 millones de parámetros, Meta implementó una estrategia económica clásica de comoditización del complemento (\textit{commoditizing the complement}). Al liberar los pesos de modelos que compiten directamente en capacidades de razonamiento y conocimiento con los sistemas cerrados de OpenAI (GPT-4o) y Google (Gemini 1.5 Pro), Meta busca destruir el poder de fijación de precios de las interfaces de programación (APIs) propietarias de sus rivales, incentivando al mismo tiempo a la comunidad global de código abierto a optimizar la pila de software sobre la cual opera su propia infraestructura publicitaria y social (Benaich \& Hogarth, 2024; Bommasani et al., 2023). La proliferación de LLaMA ha impulsado miles de desarrollos derivados mediante técnicas de ajuste fino eficiente (LoRA, QLoRA) y cuantización local, democratizando el acceso a modelos de frontera en laboratorios y empresas de todo el mundo.

\subsection{La Unión Europea y el Constitucionalismo Digital: AI Act, DMA y Efecto Bruselas}

\subsubsection{Fundamentos del Constitucionalismo Digital Europeo y Artículos 16 y 114 del TFUE}

A diferencia de la primacía del mercado en Estados Unidos o del dirigismo estatista en China, la aproximación de la Unión Europea a la inteligencia artificial emana de una doctrina jurídica consolidada: el \textbf{Constitucionalismo Digital} (\textit{Digital Constitutionalism}). Como teoriza De Gregorio (2022) y profundiza Hulok (2025), esta doctrina postula que el ciberespacio y las tecnologías algorítmicas no son territorios anárquicos exentos de soberanía, sino espacios públicos transnacionales donde el poder de las corporaciones privadas debe ser delimitado normativamente para proteger la dignidad humana, la autonomía personal y los principios del Estado democrático de derecho.

La base jurídica formal para la legislación europea sobre IA se asienta en dos pilares del Tratado de Funcionamiento de la Unión Europea (TFUE):
\begin{itemize}
    \item \textbf{El Artículo 114 del TFUE:} Faculta a las instituciones comunitarias para adoptar medidas relativas a la armonización y el correcto funcionamiento del mercado interior europeo, impidiendo que los Estados miembros fragmenten el comercio mediante regulaciones nacionales divergentes.
    \item \textbf{El Artículo 16 del TFUE:} Consagra el derecho fundamental a la protección de los datos de carácter personal y mandata el establecimiento de normas sobre el tratamiento y la libre circulación de dichos datos.
\end{itemize}

El constitucionalismo digital europeo trasciende la visión puramente procedimental de la privacidad, incorporando el principio de equidad algorítmica (\textit{algorithmic fairness}), el derecho a no ser sometido a decisiones individuales automatizadas sin supervisión humana significativa (Artículo 22 del GDPR) y la obligación de transparencia auditable sobre los procesos de inferencia estadística (Floridi, 2023; Malgieri, 2022).

\subsubsection{Arquitectura Cuatripartita de Riesgo del Reglamento de Inteligencia Artificial (EU AI Act)}

Aprobado en 2024 como el Reglamento (UE) 2024/1689, el \textbf{Reglamento de Inteligencia Artificial} (\textit{EU Artificial Intelligence Act}) constituye el primer marco jurídico exhaustivo y vinculante del mundo diseñado para regular sistemáticamente el ciclo de vida de los sistemas de IA (European Parliament \& Council of the EU, 2024). El núcleo arquitectónico de la normativa reposa sobre una clasificación piramidal graduada en función del nivel de riesgo que el sistema proyecta sobre la salud, la seguridad y los derechos fundamentales de los ciudadanos (Figura \ref{fig:piramide_riesgo}).

% ==============================================================================
% FIGURA 8.2 (TIKZ): PIRÁMIDE DE RIESGO DEL EU AI ACT Y MODELOS GPAI
% ==============================================================================
\begin{figure}[htpb!]
\centering
\begin{tikzpicture}[
    scale=0.95, transform shape,
    tier/.style={trapezium, trapezium left angle=70, trapezium right angle=70, draw=unmsm_red, line width=0.8pt, align=center, font=\small\sffamily},
    sidebox/.style={rectangle, draw=unmsm_darkred, fill=bg_marfil, line width=0.6pt, rounded corners=3pt, minimum width=4.0cm, align=left, font=\scriptsize\sffamily}
]

% Nivel 4: Riesgo Inaceptable (Cúspide)
\node[tier, fill=unmsm_red, text=white, minimum width=5.5cm, minimum height=1.1cm] (N4) at (0, 4.2) {
    \textbf{RIESGO INACEPTABLE (ARTÍCULO 5)}\\
    \scriptsize Prohibición total absoluta en territorio de la UE
};

% Nivel 3: Alto Riesgo
\node[tier, fill=unmsm_darkred!75, text=white, minimum width=8.5cm, minimum height=1.2cm] (N3) at (0, 2.9) {
    \textbf{ALTO RIESGO (ARTÍCULO 6 Y ANEXO III)}\\
    \scriptsize Evaluación de conformidad, gestión de riesgos, trazabilidad y gobernanza de datos
};

% Nivel 2: Riesgo Limitado
\node[tier, fill=bg_arena!90, text=black, minimum width=11.5cm, minimum height=1.2cm] (N2) at (0, 1.6) {
    \textbf{RIESGO LIMITADO (ARTÍCULO 50)}\\
    \scriptsize Obligaciones de transparencia, notificación al usuario y etiquetado de deepfakes
};

% Nivel 1: Riesgo Mínimo (Base)
\node[tier, fill=bg_marfil, text=black, minimum width=14.5cm, minimum height=1.2cm] (N1) at (0, 0.3) {
    \textbf{RIESGO MÍNIMO O NULO}\\
    \scriptsize Libre despliegue sin restricciones regulatorias adicionales (Filtros de spam, videojuegos)
};

% Bloque transversal GPAI
\node[rectangle, draw=unmsm_red, line width=1pt, fill=bg_arena, rounded corners=4pt, minimum width=14.5cm, minimum height=1.1cm, align=center] (GPAI) at (0, -1.2) {
    \textbf{RÉGIMEN TRANSVERSAL: MODELOS DE IA DE PROPÓSITO GENERAL (GPAI, ARTS. 51--55)}\\
    \scriptsize Umbral de Riesgo Sistémico: Cómputo acumulado $> 10^{25}$ FLOPs $\cdot$ Supervisión directa por la Oficina Europea de IA (\textit{EU AI Office})
};

% Anotaciones explicativas laterales
\node[sidebox] at (6.5, 4.2) {
    \textbf{Prácticas Prohibidas:}\\
    $\bullet$ Scoring social estatal/privado.\\
    $\bullet$ Manipulación cognitiva subliminal.\\
    $\bullet$ Categorización biométrica sensible.\\
    $\bullet$ Scraping masivo facial (CCTV).
};

\node[sidebox] at (6.5, 2.9) {
    \textbf{Sectores Anexo III:}\\
    $\bullet$ Empleo y selección de personal.\\
    $\bullet$ Evaluación crediticia y seguros.\\
    $\bullet$ Infraestructuras críticas y salud.\\
    $\bullet$ Educación y justicia penal.
};

\end{tikzpicture}
\caption{\textit{Pirámide de clasificación de riesgos del EU AI Act y régimen transversal para modelos GPAI}}
\vspace{0.1cm}
{\scriptsize \textbf{Fuente:} Elaboración propia basada en el Reglamento (UE) 2024/1689 y los análisis de Floridi (2023) y Hulok (2025).}
\label{fig:piramide_riesgo}
\end{figure}

\vspace{0.3cm}

Los cuatro niveles de la pirámide de riesgo operan de la siguiente manera:

\begin{enumerate}
    \item \textbf{Riesgo Inaceptable (Artículo 5):} Abarca prácticas algorítmicas que vulneran intrínsecamente la dignidad humana y los derechos fundamentales, quedando terminantemente prohibidas en la Unión. Incluye: (a) sistemas de puntuación social (\textit{social scoring}) administrados por entidades públicas o privadas; (b) técnicas subliminales manipulativas que alteren el comportamiento de una persona causándole daño físico o psicológico; (c) explotación intencional de vulnerabilidades asociadas a la edad, discapacidad o condición socioeconómica; (d) categorización biométrica para inferir opiniones políticas, orientación sexual, raza o convicciones religiosas; y (e) raspado indiscriminado de imágenes faciales en internet o circuitos cerrados de televisión para bases de datos de reconocimiento facial.
    \item \textbf{Alto Riesgo (Artículos 6 y Anexo III):} Comprende sistemas que no están prohibidos pero cuyo potencial de daño requiere un régimen mandatorio de aseguramiento de calidad antes de su entrada al mercado. Se clasifican aquí las aplicaciones en infraestructuras críticas (transporte, redes eléctricas), educación y formación profesional (evaluación de admisiones y exámenes), empleo y recursos humanos (cribado de currículos y asignación de tareas), servicios públicos esenciales (evaluación de solvencia crediticia y seguros médicos), gestión de fronteras y administración de justicia. Estos sistemas deben someterse a una evaluación de conformidad obligatoria, implementar sistemas formales de gestión de riesgos a lo largo de todo su ciclo de vida, garantizar la representatividad y mitigación de sesgos en los conjuntos de datos, registrar logs de auditoría técnica continua y garantizar supervisión humana efectiva (\textit{human-in-the-loop}).
    \item \textbf{Riesgo Limitado (Artículo 50):} Exige el cumplimiento de obligaciones estrictas de transparencia e información al usuario. Los proveedores deben garantizar que los usuarios sean notificados de forma explícita de que están interactuando con un agente de inteligencia artificial (chatbots conversacionales) y los contenidos sintéticos de audio, video o imagen que constituyan ultrafalsificaciones (\textit{deepfakes}) deben portar marcas de agua indelebles y metadatos legibles por máquina.
    \item \textbf{Riesgo Mínimo o Nulo:} Comprende la gran mayoría de los sistemas de software convencionales que integran funciones de IA (filtros de correo no deseado, motores de recomendación de comercio electrónico, videojuegos). Estos sistemas continúan operando sin barreras regulatorias adicionales, permitiendo la innovación desregulada en aplicaciones cotidianas.
\end{enumerate}

\subsubsection{Gobernanza de Modelos GPAI, Umbrales Sistémicos y la Oficina Europea de IA}

Durante la etapa final de deliberación legislativa del AI Act en 2023, la irrupción global de ChatGPT y los modelos fundacionales multimodales obligó a los colegisladores europeos a rediseñar la norma para evitar que los modelos base eludieran las obligaciones aplicables a las aplicaciones finales. El resultado fue la incorporación de un régimen transversal para \textbf{Modelos de IA de Propósito General} (General-Purpose AI Models, GPAI), codificado en los Artículos 51 a 55 del Reglamento (Hulok, 2025; Veale \& Borgesius, 2021).

El marco normativo distingue dos categorías operativas de modelos GPAI:
\begin{enumerate}[label=\textbf{\alph*.}]
    \item \textbf{Proveedores de Modelos GPAI Estándar:} Deben elaborar y mantener documentación técnica exhaustiva para facilitar la integración por parte de desarrolladores downstream, divulgar resúmenes suficientemente detallados del contenido utilizado para el pre-entrenamiento con respeto al derecho de autor de la UE, e implementar políticas formales de cumplimiento de la propiedad intelectual.
    \item \textbf{Modelos GPAI con Riesgo Sistémico (Artículo 51):} Se presume de manera automática que un modelo presenta riesgo sistémico cuando la cantidad acumulada de cálculo utilizada para su entrenamiento, medida en operaciones de punto flotante, supera el umbral matemático de $10^{25}$ FLOPs:
    \begin{equation}
        \text{FLOPs}_{\text{acumulado}} > 10^{25} \quad \implies \quad \text{Clasificación de Riesgo Sistémico Automática}
    \end{equation}
    Este umbral técnico sitúa inmediatamente a sistemas como GPT-4, Claude 3.5 Sonnet, Gemini 1.5 y LLaMA 3.1 bajo la supervisión de máxima intensidad. Los desarrolladores de estos modelos deben ejecutar evaluaciones formales de modelos (incluyendo simulaciones de ciberataques y fugas de control), evaluar y mitigar riesgos sistémicos a escala de la Unión, notificar incidentes de seguridad críticos a la Comisión Europea y asegurar estándares reforzados de ciberseguridad física y lógica para la infraestructura del modelo.
\end{enumerate}

Para implementar este régimen, el Reglamento creó la \textbf{Oficina Europea de IA} (\textit{EU AI Office}) dentro de la Comisión Europea. Este organismo centralizado opera como autoridad técnica de auditoría, asistido por un Panel Científico de expertos independientes con facultades para inspeccionar la arquitectura interna de los modelos, exigir la presentación de logs de entrenamiento e imponer sanciones financieras proporcionales que pueden alcanzar hasta €35.000.000 de euros o el 7\% de la facturación global anual consolidada de la empresa infractora, el monto que resulte superior (Artículo 99 del AI Act).

\subsubsection{El Marco Jurídico Conexo: DMA, GDPR, DSA y la Ley de Datos}

El AI Act no opera en el vacío institucional; se inserta dentro de una compleja red regulatoria comunitaria que redefine la economía digital europea:
\begin{itemize}
    \item \textbf{Ley de Mercados Digitales (Digital Markets Act, DMA):} Diseñada para contener el abuso de poder económico de las plataformas designadas como "guardianes de acceso" (\textit{gatekeepers}). Prohíbe prácticas de auto-preferencia (\textit{self-preferencing}), donde un guardián de acceso favorece sus propios servicios de IA sobre los de competidores que operan en su plataforma; prohíbe el empaquetamiento forzado (\textit{tying}) de modelos propietarios en sistemas operativos sin opción de desinstalación; y exige la interoperabilidad de interfaces de mensajería y APIs de datos (Zheng, 2025).
    \item \textbf{Reglamento General de Protección de Datos (GDPR):} La jurisprudencia consolidada por el Tribunal de Justicia de la Unión Europea (TJUE) y las autoridades nacionales de control establece que el raspado web masivo e indiscriminado de datos personales para alimentar corpus de pre-entrenamiento de LLMs constituye una violación del principio de licitud, lealtad y limitación de la finalidad (Artículo 5 y 6 del GDPR), careciendo de base legal legítima si los titulares de los datos no otorgaron su consentimiento explícito o si no se garantiza el derecho de supresión y oposición (Malgieri, 2022).
    \item \textbf{Ley de Servicios Digitales (DSA) y Ley de Datos (Data Act):} Imponen obligaciones de moderación de contenidos ilícitos, auditoría de algoritmos de recomendación en plataformas de muy gran escala (VLOPs) y garantizan el acceso e intercambio de datos industriales no personales generados por dispositivos conectados (IoT), impidiendo el bloqueo artificial de datos por parte de fabricantes de hardware.
\end{itemize}

\subsubsection{Evidencia Empírica del Efecto Bruselas y Fricciones Operativas Transatlánticas}

La capacidad de la Unión Europea para moldear el comportamiento de las corporaciones globales se conceptualiza a través de la teoría del \textbf{"Efecto Bruselas"} (\textit{The Brussels Effect}), formulada por la jurista Anu Bradford (2020). Bradford demuestra que, debido al tamaño del mercado interno europeo (más de 450 millones de consumidores de alto poder adquisitivo) y la indivisibilidad técnica de los productos digitales globales, las corporaciones multinacionales frecuentemente optan por adaptar sus productos y estándares de privacidad a nivel mundial a las exigencias regulatorias europeas, en lugar de mantener líneas de producción o arquitecturas de software fragmentadas para diferentes regiones.

Sin embargo, en el ámbito de la inteligencia artificial de frontera, la investigación empírica conducida por Kshetri (2024) documenta un fenómeno divergente: en lugar de una adopción global pacífica inmediata, el rigor regulatorio europeo ha provocado \textbf{fricciones operativas transatlánticas}, caracterizadas por retrasos estratégicos en los despliegues comerciales, retención deliberada de funciones y suspensiones cautelares sobre los gigantes tecnológicos estadounidenses:

\begin{itemize}
    \item \textbf{OpenAI y el Precedente Italiano del Garante (Marzo de 2023):} La Autoridad Italiana de Protección de Datos (\textit{Garante per la protezione dei dati personali}) ordenó el bloqueo inmediato y la suspensión temporal del procesamiento de datos de usuarios italianos por parte de ChatGPT. La autoridad fundamentó su intervención en la falta de base jurídica para la recolección masiva de datos personales para el entrenamiento del modelo, la ausencia total de filtros eficaces de verificación de edad para proteger a menores y la emisión de información inexacta sobre personas físicas (alucinaciones lesivas del honor). Para restablecer el servicio, OpenAI debió implementar controles de edad, formularios de exclusión (\textit{opt-out}) accesibles y clarificar sus avisos de privacidad (Kshetri, 2024).
    \item \textbf{Meta Platforms y la Retención de LLaMA Multimodal (2024):} Ante las advertencias de la Comisión de Protección de Datos de Irlanda (DPC) sobre el uso no consentido de publicaciones públicas de usuarios de Instagram y Facebook para alimentar sus modelos multimodales, Meta tomó la decisión estratégica de retrasar y retener el lanzamiento de sus funciones de IA multimodal en la Unión Europea, argumentando que la rigidez regulatoria impedía ofrecer una experiencia homogénea a los consumidores europeos.
    \item \textbf{Alphabet / Google y los Retrasos de Bard y Gemini:} El despliegue de los servicios conversacionales de Google enfrentó meses de postergación en la UE en comparación con más de 100 países de América, Asia y Oceanía, debido a requerimientos de evaluación de impacto relativos a la protección de datos (DPIA) impuestos por los reguladores comunitarios.
    \item \textbf{Anthropic y el Lanzamiento Diferido de Claude:} La startup estadounidense postergó casi doce meses el acceso comercial directo a su familia de modelos Claude en territorio europeo, priorizando mercados con menor escrutinio sobre el uso de datos de entrenamiento y derechos de autor.
\end{itemize}

Esta evidencia demuestra que las Big Tech utilizan la retención del acceso a innovaciones como mecanismo de presión política contra los reguladores comunitarios, tensionando el alcance del Efecto Bruselas tradicional.

\subsubsection{Soberanía Tecnológica, Supercómputo Público (EuroHPC) y Campeones Nativos}

Para contrarrestar el riesgo de dependencia tecnológica estructural respecto a la infraestructura estadounidense y china, la Unión Europea ha impulsado una estrategia proactiva de soberanía computacional. El buque insignia de esta iniciativa es la \textbf{Empresa Común de Informática de Alto Rendimiento Europea} (EuroHPC JU), dotada con un paquete de inversión superior a €4.000 millones de euros para el período 2024--2027 (Hulok, 2025; Kshetri, 2024).

La red EuroHPC ha desplegado una constelación de supercomputadores públicos de clase mundial, entre los que destacan:
\begin{enumerate}
    \item \textbf{MareNostrum 5 (Barcelona Supercomputing Center, España):} Sistema pre-exaescala con una partición acelerada compuesta por miles de GPUs Nvidia H100 y procesadores de arquitectura abierta, optimizado para bioinformática, modelado climático y entrenamiento de modelos fundacionales lingüísticos.
    \item \textbf{LUMI (Kajaani, Finlandia):} Supercomputador exaescala líder en eficiencia ecológica, operando con energía 100\% hidroeléctrica y reutilizando el calor residual para la red de calefacción urbana de la ciudad.
    \item \textbf{LEONARDO (CINECA, Italia):} Sistema de computación de alto rendimiento dedicado a investigación académica, física computacional y modelado algorítmico europeo.
\end{enumerate}

Complementariamente, la Comisión Europea lanzó en 2024 la iniciativa de las \textbf{"Fábricas de IA"} (\textit{AI Factories}), una red de centros dedicados que combinan tiempo de cómputo en supercomputadores EuroHPC, espacios comunes europeos de datos federados e instalaciones de soporte técnico dirigidos exclusivamente a startups e investigadores comunitarios. Este ecosistema ha permitido el rápido florecimiento de campeones tecnológicos nativos soberanos, encabezados por la empresa francesa \textbf{Mistral AI} (fundada por exinvestigadores de Meta y Google DeepMind, valorada en más de €6.000 millones de euros con su arquitectura Mixtral) y la alemana \textbf{Aleph Alpha}, orientada a la soberanía de datos y modelos con total explicabilidad y cumplimiento normativo para la administración pública y la industria automotriz.

\subsection{El Ecosistema Chino de IA: Soberanía, Guerra de Semiconductores y Modelos de Frontera}

\subsubsection{La Transición Estratégica: De Ensamblador Global a Soberanía Algorítmica}

El desarrollo de la inteligencia artificial en la República Popular China no responde a una dinámica de libre mercado espontánea, sino a un proceso de \textbf{modernización tecno-estratégica dirigida por el Estado}. El punto de inflexión fundacional ocurrió en julio de 2017 con la publicación del \textit{"Plan de Desarrollo de Inteligencia Artificial de Nueva Generación"} (新一代人工智能发展规划) por parte del Consejo de Estado. Dicho documento fijó metas cronológicas explícitas: lograr la paridad con las potencias occidentales para 2020, consolidar avances de ruptura en teorías algorítmicas fundamentales para 2025, y erigir a China como el centro de innovación global indiscutido de la inteligencia artificial hacia el año 2030 (Thomaz et al., 2022; Zheng, 2025).

Históricamente, la ventaja competitiva china descansaba en la masividad de su mercado interno (más de 1.000 millones de usuarios digitales), la disponibilidad de inmensos volúmenes de datos móviles y una industria manufacturera orientada al ensamblaje de componentes electrónicos. No obstante, las sanciones impuestas por Washington a corporaciones como ZTE y Huawei a partir de 2018 desnudaron la extrema vulnerabilidad de la economía china ante la dependencia de propiedad intelectual de semiconductores extranjeros. A partir de ese momento, el liderazgo político chino categorizó la "independencia y autosuficiencia en ciencia y tecnología" (\textit{Science and Technology Self-Reliance}) como el imperativo máximo de seguridad nacional, impulsando la sustitución integral de importaciones en toda la cadena de suministro digital.

\subsubsection{El Embargo Tecnológico del BIS, Litografía DUV/EUV y la Sustitución Nacional (Ascend 910B)}

La respuesta estratégica de Estados Unidos para frenar el ascenso tecnológico chino se articuló formalmente a través de la Oficina de Industria y Seguridad (BIS) del Departamento de Comercio. Mediante una serie de regulaciones extraterritoriales emitidas en octubre de 2022 y ampliadas en octubre de 2023 y diciembre de 2024, Washington impuso un cerco tecnológico multidimensional:
\begin{enumerate}
    \item \textbf{Restricción a Procesadores de Alto Rendimiento:} Se prohibió la venta y envío a territorio chino de aceleradores con tasas de transferencia de interconexión bidireccional superiores a 600 GB/s o capacidades de procesamiento superiores a determinados umbrales de densidad de cálculo. Esta medida bloqueó la exportación de las GPUs Nvidia A100 y H100, e invalidó posteriormente los chips recortados que Nvidia diseñó para cumplir los límites iniciales (A800, H800 y B20).
    \item \textbf{Bloqueo de Equipos de Litografía y Fabricación:} Presión diplomática y regulatoria sobre el gobierno neerlandés y japonés para revocar las licencias de exportación de ASML respecto a sus escáneres de litografía de luz ultravioleta extrema (EUV) y limitar las versiones avanzadas de ultravioleta profunda (DUV por inmersión), impidiendo el acceso a herramientas esenciales para fabricar obleas por debajo de los 14 nanómetros.
    \item \textbf{Restricción al Talento y Servicios:} Prohibición a ciudadanos, residentes permanentes y entidades jurídicas estadounidenses de prestar servicios de soporte, mantenimiento o consultoría a fundiciones de semiconductores chinas que operen nodos avanzados.
\end{enumerate}

Lejos de paralizar el desarrollo chino, el embargo aceleró drásticamente la inversión y maduración del ecosistema nacional de silicio. La fundición china SMIC (Semiconductor Manufacturing International Corporation) logró la fabricación de semiconductores de clase 7 nanómetros (nodo N+2) mediante la técnica de patrones múltiples (\textit{multi-patterning}) sobre maquinaria DUV preexistente.

En el ámbito de los aceleradores de IA, el fabricante líder nacional \textbf{Huawei} consolidó su procesador \textbf{Ascend 910B} y la posterior serie \textbf{Ascend 910C}. Diseñado con microarquitectura DaVinci v2, el Ascend 910B alcanza un rendimiento de cálculo denso en punto flotante de media precisión (FP16/BF16) comparable al 70--80\% de una GPU Nvidia A100/H100 en tareas de inferencia y entrenamiento estándar. Su ecosistema se sustenta en la pila de software de código abierto \textbf{CANN} (Compute Architecture for Neural Networks), la cual traduce operaciones de marcos como PyTorch y MindSpore hacia las unidades de procesamiento neuronal (NPUs) de Huawei, neutralizando parcialmente el foso defensivo de CUDA (Allen, 2023; Zheng, 2025).

\subsubsection{Infraestructura Nacional Federada: La Red Datos del Este, Cómputo del Oeste (Dongshu Xisuan)}

Un obstáculo físico insalvable para cualquier infraestructura de cómputo masiva es el suministro eléctrico y la disipación térmica. Mientras que en Estados Unidos los megacentros de datos se concentran en campus geográficamente restringidos (Virginia del Norte, Ohio, Texas), China implementó una solución topológica continental de naturaleza federada: la iniciativa nacional \textbf{"Datos del Este, Cómputo del Oeste"} (\textit{Dongshu Xisuan}, 东数西算), lanzada formalmente por la Comisión Nacional de Desarrollo y Reforma (NDRC) en 2022 (Zhu et al., 2025).

% ==============================================================================
% FIGURA 8.3 (TIKZ): TOPOLOGÍA DONGSHU XISUAN VS STARGATE
% ==============================================================================
\begin{figure}[htpb!]
\centering
\begin{tikzpicture}[
    scale=0.92, transform shape,
    boxStargate/.style={rectangle, draw=unmsm_darkred, fill=bg_arena, line width=0.8pt, rounded corners=3pt, minimum width=6.5cm, align=left, font=\scriptsize\sffamily},
    boxDongshu/.style={rectangle, draw=unmsm_red, fill=bg_marfil, line width=0.8pt, rounded corners=3pt, minimum width=6.5cm, align=left, font=\scriptsize\sffamily},
    header/.style={rectangle, draw=unmsm_red, fill=unmsm_red, text=white, line width=0.8pt, rounded corners=2pt, minimum width=6.5cm, align=center, font=\small\sffamily\bfseries}
]

% Lado Izquierdo: Modelo Stargate
\node[header] (H1) at (-3.8, 3.8) {MODELO OCCIDENTAL: STARGATE};
\node[boxStargate] (B1) at (-3.8, 1.8) {
    \textbf{Topología: Monolítica Hiper-Centralizada}\\
    $\bullet$ Campus masivo único o clusters coubicados.\\
    $\bullet$ Demanda eléctrica colosal concentrada: 1 a 5 GW.\\
    $\bullet$ Reactivación de centrales nucleares dedicadas.\\
    $\bullet$ Red interna InfiniBand NDR de ultra-baja latencia.\\
    $\bullet$ Dependencia crítica de la red eléctrica local.\\
    $\bullet$ Riesgo de cuellos de botella de distribución regional.
};
\node[draw=unmsm_darkred, fill=white, rounded corners=2pt, align=center, font=\scriptsize\sffamily] at (-3.8, -0.4) {
    \textbf{Ventaja:} Mínima latencia de sincronización\\
    \textbf{Límite:} Saturación térmica y disponibilidad de red
};

% Lado Derecho: Modelo Dongshu Xisuan
\node[header] (H2) at (3.8, 3.8) {MODELO CHINO: DONGSHU XISUAN};
\node[boxDongshu] (B2) at (3.8, 1.8) {
    \textbf{Topología: Red Federada Subcontinental}\\
    $\bullet$ 8 Centros de Computación y 10 Clusters Nacionales.\\
    $\bullet$ Datos generados en el Este (Shanghái, Shenzhen, Pekín).\\
    $\bullet$ Cómputo distribuido en el Oeste (Guizhou, Ningxia, Gansu).\\
    $\bullet$ Energía verde abundante: Hidroeléctrica, solar y eólica.\\
    $\bullet$ Red troncal de fibra óptica de latencia < 20 ms (OXC).\\
    $\bullet$ Optimización masiva de PUE ($< 1.2$) por clima frío.
};
\node[draw=unmsm_red, fill=white, rounded corners=2pt, align=center, font=\scriptsize\sffamily] at (3.8, -0.4) {
    \textbf{Ventaja:} Sostenibilidad energética y resiliencia\\
    \textbf{Límite:} Requiere arquitecturas tolerantes a latencia
};

% Comparador central
\draw[dashed, line width=1pt, draw=unmsm_gold] (0, 4.3) -- (0, -1.0);
\node[draw=unmsm_red, fill=bg_arena, circle, inner sep=2pt, font=\scriptsize\bfseries\color{unmsm_red}] at (0, 1.8) {VS};

\end{tikzpicture}
\caption{\textit{Contraste topológico de infraestructura computacional: Megacluster centralizado frente a Red federada nacional}}
\vspace{0.1cm}
{\scriptsize \textbf{Fuente:} Elaboración propia basada en Moussa et al. (2025) y la documentación de la Comisión Nacional de Desarrollo y Reforma de China (NDRC).}
\label{fig:topologia_computo}
\end{figure}

\vspace{0.3cm}

Como se conceptualiza en la Figura \ref{fig:topologia_computo}, el proyecto \textit{Dongshu Xisuan} explota las asimetrías geográficas de China: las megaciudades y centros industriales costeros del este (Pekín, Shanghái, la Gran Bahía de Guangdong-Hong Kong-Macao) concentran la actividad económica y la generación de datos, pero sufren de altos costos de suelo y severo estrés eléctrico. Por su parte, las provincias occidentales (Guizhou, Ningxia, Gansu, Mongolia Interior) cuentan con climas fríos que favorecen la refrigeración natural y poseen inmensas reservas energéticas renovables (hidroeléctrica, eólica y solar).

Mediante el despliegue de redes troncales de fibra óptica de ultra-alta capacidad con conmutación óptica pura (*Optical Cross-Connect*, OXC) y latencias entre centros inferiores a 20 milisegundos, el sistema canaliza las cargas de cómputo no sensibles a latencia instantánea ---especialmente las fases de pre-entrenamiento de LLMs, procesamiento por lotes y renderizado analítico--- hacia las provincias del oeste, reservando la infraestructura local del este para la inferencia en tiempo real de baja latencia. Esta arquitectura minimiza el consumo de la red energética metropolitana y permite a los centros de datos alcanzar índices de eficacia en el uso de la energía (\textit{Power Usage Effectiveness}, PUE) inferiores a 1,2, satisfaciendo estrictos estándares ambientales de descarbonización.

\subsubsection{Los Cuatro Caminos de Innovación Algorítmica China y Alibaba Qwen}

Frente a la imposibilidad de acceder en cantidades masivas a las GPUs más potentes de Silicon Valley, los laboratorios chinos debieron reemplazar la "fuerza bruta computacional" por innovaciones matemáticas en eficiencia arquitectónica. La evidencia extraída de los avances de los grupos de investigación de la UNMSM y los reportes de frontera identifica cuatro vías paradigmáticas de especialización técnica, a las que se suma la influencia global del ecosistema Qwen de Alibaba:

\begin{enumerate}
    \item \textbf{DeepSeek (DeepSeek-AI) --- La Eficiencia Algorítmica y el Razonamiento Puro:} DeepSeek conmocionó el ecosistema global de IA al demostrar que era posible entrenar y servir modelos con capacidades equivalentes a los sistemas más avanzados de OpenAI a una fracción ínfima de su costo de cómputo y consumo de memoria. Sus pilares matemáticos son tres:
    \begin{itemize}
        \item \textit{Atención Latente Multi-Cabeza (Multi-head Latent Attention, MLA):} Como se demostrará formalmente en el Caso Práctico 2, la atención estándar (MHA) genera un cuello de botella masivo en la memoria de la GPU durante la generación de texto largo debido al almacenamiento de las claves y valores en el caché (KV-Cache). DeepSeek comprime los vectores de clave y valor en un espacio latente de baja dimensión ($d_c = 512$), reduciendo el tamaño del KV-Cache en más del 93\% sin degradar la precisión del modelo.
        \item \textit{Mezcla de Expertos de Grano Fino (Fine-Grained MoE):} En lugar de utilizar pocos expertos masivos (p. ej., 8 expertos activando 2 como Mixtral), DeepSeek-V3 segmenta la capa de avance (\textit{feed-forward}) en 256 expertos especializados diminutos más un experto compartido, activando únicamente 8 expertos por token (37.000 millones de parámetros activos de un total de 671.000 millones), optimizando la utilización del hardware.
        \item \textit{DeepSeek-R1 y Aprendizaje por Refuerzo Puro (Pure RL):} A diferencia de los modelos occidentales alineados mediante miles de horas de costosa supervisión humana (RLHF), DeepSeek-R1 demostró que aplicar aprendizaje por refuerzo a gran escala directamente sobre el modelo base (\textit{DeepSeek-R1-Zero}) induce de manera espontánea cadenas de razonamiento complejas (\textit{Chain-of-Thought}), permitiendo al modelo corregir sus propios errores y verificar hipótesis matemáticas paso a paso antes de emitir una respuesta final, exhibiendo además una estricta honestidad epistémica (rechazo a alucinar ante vacíos de datos) (DeepSeek-AI, 2025).
    \end{itemize}
    \item \textbf{Kimi (Moonshot AI) --- El Dominio del Contexto Ultra-Largo:} Moonshot AI especializó su arquitectura en la eliminación de la fragmentación documental, implementando una ventana de contexto de procesamiento continuo de 2 a 5 millones de tokens. Mediante algoritmos de interpolación posicional dinámica y optimización de memoria linealizada, Kimi permite a investigadores sintetizar monografías de cientos de páginas, códigos fuente completos y archivos históricos en una única pasada computacional sin pérdida de resolución semántica (Moonshot AI, 2024).
    \item \textbf{ERNIE Bot / Wenxin Yiyan (Baidu) --- Integración con Búsqueda y Grafos de Conocimiento:} El gigante de las búsquedas Baidu estructuró su modelo ERNIE 4.0 mediante una arquitectura simbiótica que enlaza la red neuronal generativa con el mayor grafo de conocimiento bilingüe de China (cientos de miles de millones de nodos de entidades y hechos verificados). Esta integración permite a ERNIE recuperar información fáctica en tiempo real para verificar respuestas, gozando de una posición dominante en el sector corporativo y los servicios de computación en la nube empresarial de China (Baidu, 2024).
    \item \textbf{Zhipu AI (GLM-4 / AutoGLM) --- Linaje de Tsinghua y Agentes de Interfaz:} Surgida del Laboratorio de Ingeniería del Conocimiento (KEG) de la Universidad de Tsinghua, Zhipu AI desarrolló la arquitectura GLM (General Language Model), basada en un pre-entrenamiento autoregresivo bidireccional que combina las ventajas de los modelos tipo BERT y GPT. Su mayor contribución radica en la vanguardia de los agentes de software autónomos mediante \textbf{AutoGLM}, un sistema diseñado para la ejecución directa de comandos en interfaces gráficas de usuario (\textit{device-use}), manipulando pantallas de teléfonos inteligentes y computadoras de escritorio mediante razonamiento multimodal (Zhipu AI, 2024).
    \item \textbf{Alibaba Qwen (Tongyi Qianwen) --- El Estandarte del Código Abierto Chino:} La división en la nube de Alibaba adoptó una política masiva de pesos abiertos con su familia Qwen 2.5 (en variantes desde 0.5B hasta 72B parámetros densos y modelos MoE de 110B). Qwen se ha consolidado en benchmarks globales como uno de los sistemas más potentes del mundo en razonamiento matemático, codificación en Python y procesamiento de lenguajes asiáticos y occidentales, rivalizando directamente con la familia LLaMA de Meta (Zheng, 2025).
\end{enumerate}

\subsubsection{El Régimen de Gobernanza de la CAC y Valores Socialistas Fundamentales}

En el plano regulatorio interno, el desarrollo tecnológico chino está estrictamente enmarcado por las normativas emitidas por la Administración del Ciberespacio de China (CAC), en coordinación con la Comisión Nacional de Desarrollo y Reforma, el Ministerio de Ciencia y Tecnología y otros ministerios. El instrumento jurídico central son las \textbf{Medidas Provisionales para la Gestión de Servicios de Inteligencia Artificial Generativa} (生成式人工智能服务管理暂行办法), en vigor desde agosto de 2023 (CAC, 2023).

El régimen de la CAC se fundamenta en un control previo (\textit{ex ante}) de licenciamiento y registro:
\begin{itemize}
    \item \textbf{Alineamiento Político e Ideológico Obligatorio:} Las medidas exigen de forma taxativa que los contenidos generados por los modelos reflejen los "Valores Socialistas Fundamentales" y no contengan elementos que inciten a la subversión del poder estatal, el derrocamiento del sistema socialista, atenten contra la unidad nacional o dañen la soberanía y la reputación del Estado.
    \item \textbf{Evaluación de Seguridad y Registro de Algoritmos:} Todo proveedor que pretenda ofrecer servicios de IA generativa al público en general debe someter el modelo a una evaluación de seguridad formal conducida por la CAC y registrar su algoritmo en el Sistema Nacional de Registro de Algoritmos de Información de Internet. La evaluación audita las fuentes de los datos de entrenamiento, los métodos de filtrado de contenido lesivo y la tasa de rechazo del sistema ante consultas sensibles.
    \item \textbf{Etiquetado y Trazabilidad:} Mandato estricto de marcar con marcas de agua digitales perceptibles e imperceptibles todo contenido sintético generado, garantizando que el usuario pueda discernir el origen artificial del texto, audio o video.
\end{itemize}

Este marco crea un entorno donde la innovación técnica convive con límites estrictos sobre la libertad temática, generando lo que auditores externos identifican como una hiperprudencia algorítmica: los modelos chinos exhiben salvaguardas extremas de denegación ante temas de política interna o controversias históricas, redirigiendo al usuario hacia fuentes oficiales del gobierno (Cosgrove \& Cachia, 2025; Thomaz et al., 2022).

\subsection{Síntesis Comparativa Multidimensional: Silicio, Cómputo, Regulación y Modelos}

\subsubsection{Matriz Comparativa Trilateral de Ecosistemas (EE. UU., UE, China)}

Para sintetizar la complejidad estructural del panorama global, la Tabla \ref{tab:matriz_trilateral} presenta una matriz comparativa multidimensional que contrasta los pilares jurídicos, tecnológicos, económicos y normativos de los tres bloques hegemónicos analizados.

% ==============================================================================
% TABLA 8.1: MATRIZ COMPARATIVA TRILATERAL
% ==============================================================================
\begin{table}[htpb!]
\centering
\caption{\textbf{Tabla 8.1}\\ \textit{Matriz comparativa trilateral de ecosistemas de inteligencia artificial}}
\label{tab:matriz_trilateral}
\vspace{0.15cm}
\footnotesize
\renewcommand{\arraystretch}{1.35}
\begin{tabularx}{\textwidth}{>{\raggedright\arraybackslash\bfseries}p{2.6cm} >{\raggedright\arraybackslash}X >{\raggedright\arraybackslash}X >{\raggedright\arraybackslash}X}
\toprule
\textbf{Dimensión} & \textbf{Estados Unidos} & \textbf{Unión Europea} & \textbf{República Popular China} \\
\midrule
\textbf{Filosofía de Gobernanza} & Primacía de la innovación de mercado, desregulación pro-competitiva y contención geopolítica externa. & Constitucionalismo digital basado en riesgos \textit{ex ante}, dignidad humana y protección de derechos fundamentales. & Dirigismo estatista planificado, autosuficiencia nacional y soberanía estricta sobre los contenidos informativos. \\
\addlinespace[0.2cm]
\textbf{Instrumentos Legales Centrales} & Orden Ejecutiva 14179 (revocando EO 14110), leyes antimonopolio Sherman/Clayton, regulaciones BIS. & Reglamento de IA (EU AI Act 2024/1689), DMA, GDPR, Ley de Servicios Digitales (DSA) y Ley de Datos. & Medidas Provisionales para la Gestión de IA Generativa (CAC, 2023) y Normas de Registro Algorítmico. \\
\addlinespace[0.2cm]
\textbf{Modelo de Infraestructura} & Megacentros de datos centralizados a hiperescala y megaproyectos privados (Proyecto Stargate \$500.000M). & Red pública de supercomputación científica (EuroHPC: MareNostrum, LUMI) y Fábricas de IA (€4.000M). & Red federada continental \textit{Dongshu Xisuan} (8 hubs, 10 megaclusters, energía hidro/solar en el oeste). \\
\addlinespace[0.2cm]
\textbf{Dominio de Silicio y Aceleradores} & Monopolio de diseño y software: Nvidia (H100/B200, CUDA), AMD (MI300X); manufactura externa vía TSMC. & Liderazgo monopólico en litografía óptica avanzada: ASML (EUV/High-NA); arquitectura de IP cores ARM. & Sustitución doméstica acelerada: Huawei Ascend (910B/910C), SMIC (nodo 7nm DUV), software CANN. \\
\addlinespace[0.2cm]
\textbf{Ecosistema de Modelos} & Dominio en modelos comerciales propietarios (OpenAI, Google, Anthropic) y pesos abiertos con Meta LLaMA. & Modelos fundacionales soberanos de alta eficiencia: Mistral AI (Mixtral 8x22B) y Aleph Alpha (explicabilidad). & Eficiencia matemática extrema y pesos abiertos: DeepSeek (MLA/MoE), Kimi (2M context), ERNIE, Zhipu, Qwen. \\
\addlinespace[0.2cm]
\textbf{Estructura de Capital} & Capital de riesgo privado masivo (Silicon Valley) e integración de flujos de caja de Big Tech. & Financiamiento mixto: fondos de capital de riesgo europeos y subvenciones directas de la Comisión Europea. & Subsidios estatales directos, fondos de orientación gubernamental e inversión de gigantes digitales (Baidu, Alibaba). \\
\addlinespace[0.2cm]
\textbf{Impacto Internacional} & Atracción global de talento e imposición de estándares tecnológicos \textit{de facto} mediante software y APIs. & El "Efecto Bruselas": regulación extraterritorial sistemática forzando cumplimiento en mercados foráneos. & Exportación de infraestructura de telecomunicaciones y tecnologías inteligentes vía la "Ruta de la Seda Digital". \\
\bottomrule
\end{tabularx}
\vspace{0.1cm}
{\scriptsize \textbf{Fuente:} Elaboración propia basada en Bradford (2020), Dylag (2025), Moussa et al. (2025), Zheng (2025) y Zhu et al. (2025).}
\end{table}

\subsubsection{Economía del Cómputo, Leyes de Escalamiento y Costos de Inferencia}

Durante el lustro comprendido entre 2019 y 2024, la investigación en modelos fundacionales estuvo dominada por la hipótesis de las \textbf{Leyes de Escalamiento} (\textit{Scaling Laws}), formalizadas empíricamente por Kaplan et al. (2020) y refinadas por Hoffmann et al. (2022) en el marco de optimización de cómputo de Chinchilla. De acuerdo con estas leyes, la pérdida de entropía cruzada en modelos de lenguaje exhibe una relación de ley de potencia con el número de parámetros del modelo ($N$), el tamaño del conjunto de datos de entrenamiento ($D$) y el presupuesto de cómputo total empleado ($C$):
\begin{equation}
    L(N, D) = E + \frac{A}{N^{\alpha}} + \frac{B}{D^{\beta}}
\end{equation}
donde $E, A, B, \alpha, \beta$ son constantes empíricas y el cómputo total de entrenamiento para arquitecturas densas se aproxima mediante la regla empírica clásica:
\begin{equation}
    C \approx 6 \times N \times D \quad [\text{FLOPs}]
\end{equation}

Esta formulación incentivó una carrera por fuerza bruta computacional en Silicon Valley, multiplicando los tamaños de los modelos y exigiendo presupuestos de entrenamiento superiores a los cientos de millones de dólares. Sin embargo, en 2024--2025 se evidenció un rendimiento decreciente del escalamiento en pre-entrenamiento, trasladando el foco de la competencia hacia la **eficiencia en inferencia** y el **cómputo en tiempo de prueba** (\textit{test-time compute}) (DeepSeek-AI, 2025).

La Tabla \ref{tab:modelos_costos} expone las especificaciones arquitectónicas, volúmenes de cómputo y costos de inferencia por millón de tokens generados, demostrando la dramática reducción de costos introducida por los modelos chinos de alta eficiencia.

% ==============================================================================
% TABLA 8.2: ESPECIFICACIONES TÉCNICAS Y COSTOS DE MODELOS
% ==============================================================================
\begin{table}[htpb!]
\centering
\caption{\textbf{Tabla 8.2}\\ \textit{Especificaciones arquitectónicas, cómputo y costos de modelos de frontera globales}}
\label{tab:modelos_costos}
\vspace{0.15cm}
\footnotesize
\renewcommand{\arraystretch}{1.3}
\begin{tabularx}{\textwidth}{>{\raggedright\arraybackslash\bfseries}p{2.4cm} >{\raggedright\arraybackslash}p{2.1cm} >{\raggedright\arraybackslash}p{2.3cm} >{\centering\arraybackslash}p{1.7cm} >{\centering\arraybackslash}p{1.4cm} >{\centering\arraybackslash}p{1.5cm} >{\raggedleft\arraybackslash}p{1.8cm}}
\toprule
\textbf{Modelo} & \textbf{Organización} & \textbf{Arquitectura Base} & \textbf{Parámetros Totales/Act.} & \textbf{Ventana Contexto} & \textbf{Cómputo (FLOPs)} & \textbf{Costo Infer. (\$/1M Tok.)} \\
\midrule
\textbf{GPT-4o} & OpenAI (EE. UU.) & Transformador Denso Multimodal & Privado ($\sim$1.8T MoE est.) & 128k & $> 10^{25}$ & \$2.50 / \$10.00 \\
\textbf{Claude 3.5 Son.} & Anthropic (EE. UU.) & Transformador Híbrido & Privado (Denso) & 200k & $> 10^{25}$ & \$3.00 / \$15.00 \\
\textbf{Gemini 1.5 Pro} & Google (EE. UU.) & MoE Multimodal Nativo & Privado & 2.000k & $> 10^{25}$ & \$1.25 / \$5.00 \\
\textbf{LLaMA 3.1 405B} & Meta (EE. UU.) & Transformador Denso GQA & 405B / 405B & 128k & $\approx 3.8 \times 10^{25}$ & \$0.80 / \$2.40 (self-host) \\
\textbf{Mistral Large 2} & Mistral AI (UE) & Transformador Denso & 123B / 123B & 128k & $\approx 1.2 \times 10^{25}$ & \$2.00 / \$6.00 \\
\textbf{DeepSeek-V3} & DeepSeek (China) & MoE Fino + MLA + FP8 & 671B / 37B & 128k & $\approx 2.8 \times 10^{24}$ & \textbf{\$0.14 / \$0.28} \\
\textbf{DeepSeek-R1} & DeepSeek (China) & Pure RL + MoE + MLA & 671B / 37B & 128k & $\approx 3.0 \times 10^{24}$ & \textbf{\$0.55 / \$2.19} \\
\textbf{Kimi 1.5} & Moonshot (China) & Atención Eficiente Contexto & Privado & \textbf{2.000k} & $\approx 5.0 \times 10^{24}$ & \$0.35 / \$1.40 \\
\textbf{ERNIE 4.0 Tur.} & Baidu (China) & Enlace Grafo Conocimiento & Privado & 128k & $\approx 4.0 \times 10^{24}$ & \$0.40 / \$1.20 \\
\textbf{Qwen 2.5 72B} & Alibaba (China) & Transformador Denso GQA & 72B / 72B & 128k & $\approx 8.5 \times 10^{24}$ & \$0.20 / \$0.40 \\
\bottomrule
\end{tabularx}
\vspace{0.1cm}
{\scriptsize \textbf{Nota:} Costos de inferencia expresados en dólares estadounidenses por millón de tokens (entrada / salida) a tarifas de API pública de 2025. Parámetros y cómputo basados en DeepSeek-AI (2025), Moonshot AI (2024) y Benaich \& Hogarth (2024).}
\end{table}

\subsubsection{Marcos Regulatorios Globales y Convergencias/Divergencias Jurídicas}

En el plano jurídico formal, la comparación sistemática de las facultades de supervisión, las sanciones financieras y los requerimientos de auditoría revela que el mundo no camina hacia un estándar único armonizado, sino hacia un arbitraje regulatorio complejo (Tabla \ref{tab:marcos_regulatorios}).

% ==============================================================================
% TABLA 8.3: MARCO COMPARADO DE REGÍMENES REGULATORIOS GLOBALES
% ==============================================================================
\begin{table}[htpb!]
\centering
\caption{\textbf{Tabla 8.3}\\ \textit{Comparación de regímenes regulatorios de inteligencia artificial}}
\label{tab:marcos_regulatorios}
\vspace{0.15cm}
\footnotesize
\renewcommand{\arraystretch}{1.3}
\begin{tabularx}{\textwidth}{>{\raggedright\arraybackslash\bfseries}p{2.3cm} >{\raggedright\arraybackslash}p{2.4cm} >{\raggedright\arraybackslash}X >{\raggedright\arraybackslash}X >{\raggedright\arraybackslash}X}
\toprule
\textbf{Régimen} & \textbf{Autoridad Líder} & \textbf{Enfoque de Supervisión} & \textbf{Sanciones Máximas} & \textbf{Obligaciones Clave} \\
\midrule
\textbf{Estados Unidos (Federal)} & FTC, DOJ, NIST, Sectoriales (SEC, FDA) & \textit{Ex post} basado en daños reales al consumidor, libre competencia y ciberseguridad nacional. & Multas por prácticas desleales y desinversión forzada de activos corporativos. & Cumplimiento voluntario de estándares NIST, reporte de vulnerabilidades biológicas en defensa. \\
\addlinespace[0.15cm]
\textbf{Unión Europea (EU AI Act)} & Oficina Europea de IA (\textit{EU AI Office}) y DPAs nacionales & \textit{Ex ante} preventivo basado en pirámide de riesgos y evaluación de conformidad vinculante. & Hasta €35.000.000 o 7\% de la facturación global anual (el monto mayor). & Prohibición de prácticas inaceptables, gestión de riesgos en alto riesgo, transparencia en GPAI. \\
\addlinespace[0.15cm]
\textbf{China (CAC)} & Administración del Ciberespacio de China (CAC) y Ministerios & \textit{Ex ante} político-ideológico mediante licenciamiento, auditoría de datos y registro de algoritmos. & Suspensión de servicio, revocación de licencias comerciales y multas administrativas. & Alineamiento con Valores Socialistas, evaluación de seguridad previa al lanzamiento, marcas de agua. \\
\bottomrule
\end{tabularx}
\vspace{0.1cm}
{\scriptsize \textbf{Fuente:} Elaboración propia basada en Dylag (2025), Hulok (2025), Zheng (2025) y la normativa oficial citada.}
\end{table}

\subsubsection{Evaluación Empírica de Desempeño y Patologías por Tarea}

Para evitar concepciones idealizadas sobre cualquiera de las tecnologías analizadas, la Tabla \ref{tab:patologias_modelos} sistematiza la evidencia empírica obtenida de pruebas de rendimiento estandarizadas, identificando las fortalezas operativas específicas junto con las limitaciones o patologías críticas que caracterizan a cada modelo en despliegues reales.

% ==============================================================================
% TABLA 8.4: EVALUACIÓN DE DESEMPEÑO Y PATOLOGÍAS CRÍTICAS
% ==============================================================================
\begin{table}[htpb!]
\centering
\caption{\textbf{Tabla 8.4}\\ \textit{Desempeño empírico y patologías críticas observadas por tarea en modelos de frontera}}
\label{tab:patologias_modelos}
\vspace{0.15cm}
\footnotesize
\renewcommand{\arraystretch}{1.35}
\begin{tabularx}{\textwidth}{>{\raggedright\arraybackslash\bfseries}p{2.6cm} >{\raggedright\arraybackslash}p{2.2cm} >{\raggedright\arraybackslash}X >{\raggedright\arraybackslash}X >{\raggedright\arraybackslash}p{2.2cm}}
\toprule
\textbf{Tarea Evaluada} & \textbf{Modelo} & \textbf{Fortaleza Observada} & \textbf{Patología o Limitación Crítica} & \textbf{Mitigación Recomendada} \\
\midrule
\textbf{Razonamiento Matemático y Código Complejo} & DeepSeek-R1 (China) & Autocorrección espontánea mediante cadena de pensamiento explícita (\textit{Chain-of-Thought}); alta precisión en pruebas formales (AIME, MATH). & Mayor latencia por generación de tokens de razonamiento oculto; sensibilidad extrema al prompt de sistema. & Ajuste fino de destilación (\textit{distilled models}) para tareas de baja latencia. \\
\addlinespace[0.15cm]
\textbf{Síntesis de Documentos Extensos (1M+ Tok)} & Kimi 1.5 (Moonshot) & Retención uniforme a lo largo de toda la ventana de contexto (\textit{Needle in a Haystack} $> 99\%$). & Pérdida marginal de detalles numéricos exactos en tablas densas distribuidas en cientos de páginas. & Extracción complementaria basada en expresiones regulares o RAG estructurado. \\
\addlinespace[0.15cm]
\textbf{Consulta General e Integración Empresarial} & ERNIE 4.0 (Baidu) & Verificación fáctica superior en mandarín y sincronización nativa con la nube y bases de datos empresariales. & Sesgo corporativo pro-servicios de Baidu; menor fluidez en estructuras gramaticales de lenguas occidentales. & Despliegue en arquitecturas híbridas con modelos bilingües complementarios. \\
\addlinespace[0.15cm]
\textbf{Agentes Autónomos en Dispositivos (UI)} & AutoGLM (Zhipu AI) & Planificación multietapa y ejecución directa de clics y gestos en interfaces táctiles de telefonía y web. & Alucinaciones agenticas (declarar falsamente haber completado una acción de compra o envío que falló). & Bucles obligatorios de confirmación humana (\textit{Human-in-the-Loop}) en acciones críticas. \\
\bottomrule
\end{tabularx}
\vspace{0.1cm}
{\scriptsize \textbf{Fuente:} Elaboración propia basada en las pruebas empíricas reportadas por Cosgrove \& Cachia (2025), DeepSeek-AI (2025), Moonshot AI (2024) y Zhipu AI (2024).}
\end{table}

\subsubsection{Fricciones Geopolíticas, Desacoplamiento y Escenarios Futuros de la IA de Frontera}

La intersección entre cómputo, energía y regulación perfila tres escenarios futuros para la próxima década:
\begin{enumerate}
    \item \textbf{Fragmentación de la Red (\textit{AI Splinternet}):} El endurecimiento recíproco de sanciones y barreras normativas podría culminar en un desacoplamiento tecnológico irreversible, donde Occidente opere exclusivamente sobre ecosistemas cerrados basados en arquitecturas de hardware de Nvidia y normas del AI Act, mientras China y países alineados del Sur Global adopten pilas soberanas basadas en silicio Ascend, redes federadas y pesos abiertos de DeepSeek y Qwen.
    \item \textbf{Guerra de Fricción Económica y Burbuja de Capital:} Las astronómicas valoraciones de las empresas de infraestructura (\$500.000 millones en proyectos como Stargate) enfrentan la amenaza de la comoditización impulsada por la eficiencia china. Si modelos que demandan un costo diez veces menor de cómputo logran paridad de razonamiento, la rentabilidad esperada del supercómputo occidental podría colapsar, desatando una crisis en los mercados financieros de capital de riesgo (Moussa et al., 2025).
    \item \textbf{Convergencia Pragmática por Interdependencia Asimétrica:} A pesar de la retórica hostil, ningún bloque posee autarquía absoluta. Occidente depende de las fundiciones asiáticas y de los pesos abiertos chinos para acelerar su propia investigación; China continúa necesitando herramientas de diseño occidentales y acceso a mercados globales para amortizar sus inversiones; y la Unión Europea requiere el poder computacional de los gigantes globales para mantener su productividad. Esta interdependencia material impone límites pragmáticos al desacoplamiento total (Zheng, 2025).
\end{enumerate}

\vspace{0.4cm}

% ==============================================================================
% BLOQUE 4: CASOS PRÁCTICOS APLICADOS
% ==============================================================================
\newpage
\subsection{Casos Prácticos Aplicados}

A continuación, se presentan dos casos prácticos aplicados que operacionalizan formalmente la tríada metodológica (\textbf{Contexto}, \textbf{Análisis e Implementación}, \textbf{Interpretación y Resultados}), integrando algoritmos de auditoría de cumplimiento regulatorio y modelos matemáticos de simulación de costos computacionales.

\subsubsection{Caso Práctico 1: Impacto Operativo del "Efecto Bruselas" y Auditoría de Cumplimiento Técnico bajo el EU AI Act}

\subsubsection*{Contexto:}
Una institución bancaria multinacional con sede operativa en Nueva York y filiales en Madrid y Fráncfort planea desplegar un sistema avanzado de inteligencia artificial generativa de propósito general (GPAI). El sistema cumplirá dos funciones simultáneas: (1) asistencia automatizada y perfilado de riesgo para la concesión de créditos hipotecarios y comerciales a clientes corporativos y particulares; y (2) cribado y selección automatizada de candidatos en los procesos de contratación del departamento de recursos humanos.

El modelo fundacional subyacente fue entrenado en Estados Unidos sobre una infraestructura de hiperescala acumulando un cómputo estimado de $3.2 \times 10^{25}$ FLOPs. La entidad bancaria enfrenta el dilema operativo documentado por Kshetri (2024): lanzar el producto de forma acelerada asumiendo los estándares de mercado estadounidenses, o someter el sistema a una auditoría formal de cumplimiento técnico bajo el Reglamento de IA de la UE (Reglamento UE 2024/1689), arriesgando sanciones que, bajo los Artículos 99 y 101, pueden alcanzar hasta €35.000.000 o el 7\% de los ingresos anuales globales por infracciones relativas a prácticas prohibidas, y hasta el 3\% o €15.000.000 por incumplimientos en sistemas de alto riesgo.

\subsubsection*{Análisis e Implementación:}
Para mitigar el riesgo de incumplimiento catastrófico y automatizar la verificación estatutaria previa al despliegue, el equipo de ingeniería y cumplimiento legal diseña un motor de auditoría computacional en Python denominado \texttt{EU\_AI\_Compliance\_Engine}. El algoritmo evalúa sistemáticamente cuatro capas de control:
\begin{enumerate}
    \item \textbf{Criba de Prácticas Prohibidas (Artículo 5):} Detección de scoring social discriminatorio o manipulación cognitiva.
    \item \textbf{Clasificación de Alto Riesgo (Artículo 6 y Anexo III):} Determinación de si el caso de uso pertenece a los sectores regulados (empleo o acceso a servicios financieros esenciales).
    \item \textbf{Evaluación de Riesgo Sistémico en GPAI (Artículo 51):} Verificación del umbral de cómputo ($> 10^{25}$ FLOPs).
    \item \textbf{Cálculo de Sanciones Potenciales y Hoja de Ruta de Mitigación:} Generación automatizada de las salvaguardas mandatorias (auditoría de sesgos en datos, logs de trazabilidad y supervisión humana).
\end{enumerate}

El Listado \ref{code:compliance_engine} detalla la implementación formal del motor de auditoría.

\begin{tcolorbox}[colback=bg_marfil, colframe=unmsm_gold, arc=3pt, boxrule=0.6pt, top=0.3cm, bottom=0.3cm, left=0.4cm, right=0.4cm]
\label{code:compliance_engine}
\scriptsize
\begin{verbatim}
# compliance_audit_engine.py
# Motor Automatizado de Auditoria de Conformidad Regulatoria (EU AI Act)
from dataclasses import dataclass
from typing import List, Dict, Tuple

@dataclass
class ModelMetadata:
    name: str
    training_compute_flops: float
    use_case_sector: str
    features: List[str]
    global_annual_turnover_eur: float

class EUAIActAuditor:
    PROHIBITED_PRACTICES = {
        "social_scoring", "subliminal_manipulation", 
        "biometric_categorization_sensitive", "indiscriminate_cctv_scraping"
    }
    
    HIGH_RISK_SECTORS = {
        "employment_hr", "credit_scoring_finance", "critical_infrastructure",
        "education_grading", "justice_law_enforcement", "biometrics"
    }
    
    SYSTEMIC_RISK_FLOP_THRESHOLD = 1e25

    def __init__(self, metadata: ModelMetadata):
        self.meta = metadata
        self.audit_log: List[str] = []
        self.risk_tier: str = "Minimal"
        self.is_systemic_gpai: bool = False
        self.max_penalty_eur: float = 0.0

    def evaluate_compliance(self) -> Dict[str, any]:
        # 1. Auditoria de Practicas Prohibidas (Articulo 5)
        detected_prohibitions = set(self.meta.features).intersection(self.PROHIBITED_PRACTICES)
        if detected_prohibitions:
            self.risk_tier = "Prohibited"
            # Sancion Articulo 99(3): Hasta 35M EUR o 7% de la facturacion global
            penalty_7pct = self.meta.global_annual_turnover_eur * 0.07
            self.max_penalty_eur = max(35_000_000.0, penalty_7pct)
            self.audit_log.append(f"VIOLACION ART. 5: Practicas prohibidas detectadas: {detected_prohibitions}")
            return self._generate_report(approved=False)

        # 2. Auditoria de Alto Riesgo (Articulo 6 y Anexo III)
        if self.meta.use_case_sector in self.HIGH_RISK_SECTORS:
            self.risk_tier = "High Risk"
            # Sancion Articulo 99(4): Hasta 15M EUR o 3% de la facturacion
            penalty_3pct = self.meta.global_annual_turnover_eur * 0.03
            self.max_penalty_eur = max(15_000_000.0, penalty_3pct)
            self.audit_log.append(f"ALTO RIESGO DETECTADO (Anexo III): Sector '{self.meta.use_case_sector}'.")

        # 3. Auditoria de Modelos GPAI con Riesgo Sistemico (Articulo 51)
        if self.meta.training_compute_flops > self.SYSTEMIC_RISK_FLOP_THRESHOLD:
            self.is_systemic_gpai = True
            self.audit_log.append(
                f"MODELO GPAI CON RIESGO SISTEMICO (Art. 51): Computo acumulado "
                f"{self.meta.training_compute_flops:.2e} FLOPs supera umbral de 1e25."
            )

        return self._generate_report(approved=True)

    def _generate_report(self, approved: bool) -> Dict[str, any]:
        requirements = []
        if self.risk_tier == "High Risk":
            requirements.extend([
                "Art. 9: Sistema de gestion de riesgos continuo.",
                "Art. 10: Gobernanza de datos de entrenamiento y mitigacion de sesgos.",
                "Art. 12: Registro automatico de logs tecnicos de operacion.",
                "Art. 14: Mecanismos de supervision humana efectiva (Human-in-the-Loop)."
            ])
        if self.is_systemic_gpai:
            requirements.extend([
                "Art. 53: Evaluaciones de vulnerabilidad y pruebas de penetracion (red-teaming).",
                "Art. 55: Notificacion obligatoria de incidentes graves a la Oficina Europea de IA."
            ])
        return {
            "model_name": self.meta.name,
            "deployment_approved": approved,
            "assigned_risk_tier": self.risk_tier,
            "systemic_gpai": self.is_systemic_gpai,
            "max_penalty_exposure_eur": self.max_penalty_eur,
            "mandatory_compliance_checklist": requirements,
            "audit_trail": self.audit_log
        }

# Instanciacion del caso practico:
bank_model = ModelMetadata(
    name="FinHR-Frontier-v1",
    training_compute_flops=3.2e25,
    use_case_sector="credit_scoring_finance",
    features=["automated_loan_approval", "candidate_cv_ranking"],
    global_annual_turnover_eur=1_200_000_000.0 # 1.200 Millones EUR
)

auditor = EUAIActAuditor(bank_model)
report = auditor.evaluate_compliance()
\end{verbatim}
\end{tcolorbox}

\subsubsection*{Interpretación y Resultados:}
La ejecución del motor de auditoría sobre el sistema bancario genera un diagnóstico determinante:
\begin{enumerate}
    \item \textbf{Clasificación de Riesgo:} El sistema queda categorizado automáticamente en \textbf{Alto Riesgo} bajo el Anexo III (servicios financieros esenciales y empleo), descartando prácticas prohibidas del Artículo 5.
    \item \textbf{Calificación de Modelo Sistémico GPAI:} Al registrar $3.2 \times 10^{25}$ FLOPs, el modelo subyacente supera el umbral del Artículo 51, quedando sujeto a la jurisdicción directa de la Oficina Europea de IA (\textit{EU AI Office}).
    \item \textbf{Exposición a Multas:} La exposición potencial a sanciones asciende a \textbf{€36.000.000 de euros} (calculada como el 3\% de la facturación global de €1.200M, que supera el piso fijo de €15M).
    \item \textbf{Impacto en Tiempo y Costos:} Adaptar el sistema para cumplir con los Artículos 9, 10, 12 y 14 exige incorporar una canalización de auditoría de sesgos demográficos (evitando discriminación por género o código postal en los créditos) y habilitar una interfaz de control donde ningún préstamo o descarte laboral sea ejecutado sin validación explícita de un oficial bancario humano. Este proceso introduce un retraso de 4 meses en el lanzamiento europeo frente a la versión de EE. UU., validando empíricamente la tesis de Kshetri (2024) sobre el impacto operativo del Efecto Bruselas.
\end{enumerate}

\subsubsection{Caso Práctico 2: Simulación Técnico-Económica de Escalamiento de Cómputo: Megacluster Hiperescalar frente a Clúster Federado de Eficiencia Algorítmica}

\subsubsection*{Contexto:}
Un consorcio científico busca entrenar y servir a escala de producción un modelo fundacional de capacidades avanzadas. El equipo de investigación debe tomar una decisión estratégica de arquitectura e infraestructura entre dos paradigmas geopolíticamente divergentes:
\begin{itemize}
    \item \textbf{Opción A (Paradigma Occidental de Escalamiento Denso):} Adquisición de tiempo en un megacluster de hiperescala equipado con 16.384 GPUs Nvidia H100 SXM5 (80GB HBM3, 700W TDP, interconexión NVLink 4 de 900 GB/s e InfiniBand NDR), entrenando una arquitectura densa clásica de 405.000 millones de parámetros con atención multi-cabeza (MHA) en precisión BF16 (arquetipo LLaMA 3.1).
    \item \textbf{Opción B (Paradigma Chino de Sustitución y Eficiencia Algorítmica):} Despliegue en un clúster federado de infraestructura nacional equipado con 20.480 aceleradores Huawei Ascend 910B (64GB HBM2e, 400W TDP, interconexión HCCS de 392 GB/s y redes RoCE v2), entrenando una arquitectura de Mezcla de Expertos de Grano Fino (MoE) de 671.000 millones de parámetros totales pero solo 37.000 millones de parámetros activos por token, con Atención Latente Multi-Cabeza (MLA) y precisión nativa FP8 (arquetipo DeepSeek-V3).
\end{itemize}

El problema radica en determinar si la eficiencia matemática del diseño algorítmico chino logra compensar el déficit en el ancho de banda de memoria física (HBM2e vs. HBM3) y potencia de interconexión impuesto por el embargo de semiconductores.

\subsubsection*{Análisis e Implementación:}
Para modelar rigurosamente la física y la economía de ambos sistemas, se formalizan las ecuaciones gobernantes de consumo de memoria y Costo Total de Propiedad (TCO):

\begin{enumerate}
    \item \textbf{Consumo de Memoria de KV-Cache en Inferencia:}
    En una arquitectura estándar de Atención Multi-Cabeza (MHA), para una longitud de contexto $s$, un tamaño de lote $b$, un número de capas $n_{\text{capas}}$, un número de cabezas $n_{\text{cabezas}}$ y una dimensión de cabeza $d_{\text{cabeza}}$, el consumo de memoria del caché por token está dado por:
    \begin{equation}
        \text{Mem}_{\text{MHA}} = 2 \times n_{\text{capas}} \times n_{\text{cabezas}} \times d_{\text{cabeza}} \times \text{bytes\_per\_elem} \quad \left[\frac{\text{Bytes}}{\text{token}}\right]
    \end{equation}
    En la Atención Latente Multi-Cabeza (MLA) de DeepSeek, los vectores de clave y valor se comprimen en un vector latente de dimensión $d_c$, junto con un vector rotatorio desacoplado de dimensión $d_R$:
    \begin{equation}
        \text{Mem}_{\text{MLA}} = 2 \times n_{\text{capas}} \times (d_c + d_R) \times \text{bytes\_per\_elem} \quad \left[\frac{\text{Bytes}}{\text{token}}\right]
    \end{equation}
    Para $n_{\text{capas}} = 61$, $n_{\text{cabezas}} \times d_{\text{cabeza}} = 128 \times 128 = 16.384$, frente a $d_c + d_R = 512 + 64 = 576$, la compresión matemática del KV-Cache es:
    \begin{equation}
        \text{Ratio de Compresión} = 1 - \frac{576}{16.384} \approx 96.48\% \text{ de reducción en memoria de contexto}
    \end{equation}

    \item \textbf{Modelo de Costo Total de Propiedad (TCO) y Costo por Token:}
    El costo total de operar el clúster durante un ciclo anual se descompone en gasto de capital amortizado (\textit{CapEx}) y gasto operativo energético (\textit{OpEx}):
    \begin{equation}
        \text{TCO} = N_{\text{chips}} \times \left( \frac{\text{Costo}_{\text{chip}} + \text{Costo}_{\text{red}}}{T_{\text{amortización}}} \right) + N_{\text{chips}} \times P_{\text{chip}} \times \text{PUE} \times 8760\,\text{h} \times \text{Tarifa}_{\text{kWh}} + \text{OpEx}_{\text{mant}}
    \end{equation}
\end{enumerate}

El Listado \ref{code:tco_simulation} presenta la simulación numérica integral en Python.

\begin{tcolorbox}[colback=bg_marfil, colframe=unmsm_gold, arc=3pt, boxrule=0.6pt, top=0.3cm, bottom=0.3cm, left=0.4cm, right=0.4cm]
\label{code:tco_simulation}
\scriptsize
\begin{verbatim}
# compute_tco_simulation.py
# Simulacion Tecnico-Economica Comparada de Supercomputo e Inferencia
from dataclasses import dataclass

@dataclass
class ClusterSpecs:
    name: str
    num_accelerators: int
    chip_cost_usd: float
    network_cost_per_chip_usd: float
    power_watts_per_chip: float
    pue: float
    electricity_cost_kwh: float
    hbm_bandwidth_tb_s: float
    active_params_billion: float
    kv_cache_dim_per_layer: int
    num_layers: int
    precision_bytes: float # 2 para BF16, 1 para FP8

def simulate_cluster_performance(c: ClusterSpecs, batch_size: int, seq_len: int):
    # 1. Calculo de Memoria KV-Cache por Token (Bytes)
    # Factor 2 para Keys y Values
    kv_cache_bytes_per_token = 2 * c.num_layers * c.kv_cache_dim_per_layer * c.precision_bytes
    total_kv_cache_gb = (kv_cache_bytes_per_token * batch_size * seq_len) / (1024**3)

    # 2. Calculo de TCO Anual (Amortizacion 3 anos)
    capex_per_chip = (c.chip_cost_usd + c.network_cost_per_chip_usd) / 3.0
    total_capex_annual = c.num_accelerators * capex_per_chip
    
    # Consumo Electrico Anual (MWh)
    total_power_kw = (c.num_accelerators * c.power_watts_per_chip * c.pue) / 1000.0
    annual_energy_kwh = total_power_kw * 8760.0
    total_opex_energy = annual_energy_kwh * c.electricity_cost_kwh
    total_opex_maintenance = total_capex_annual * 0.10 # 10% mantenimiento
    
    tco_annual_usd = total_capex_annual + total_opex_energy + total_opex_maintenance

    # 3. Rendimiento de Inferencia Estimado (Tokens por Segundo Saturado por Ancho de Banda)
    # Memoria leida por token = parametros_activos * precision
    bytes_per_token_weights = c.active_params_billion * 1e9 * c.precision_bytes
    cluster_memory_bw_tb_s = c.num_accelerators * c.hbm_bandwidth_tb_s
    max_tokens_per_sec = (cluster_memory_bw_tb_s * 1e12) / bytes_per_token_weights

    # Costo por Millon de Tokens Generados ($/1M tokens)
    annual_tokens_generated = max_tokens_per_sec * 8760 * 3600 * 0.40 # 40% utilizacion
    cost_per_million_tokens = (tco_annual_usd / (annual_tokens_generated / 1e6))

    return {
        "cluster": c.name,
        "kv_cache_gb_per_batch": round(total_kv_cache_gb, 2),
        "tco_annual_million_usd": round(tco_annual_usd / 1e6, 2),
        "annual_energy_gwh": round(annual_energy_kwh / 1e6, 2),
        "throughput_tokens_sec": round(max_tokens_per_sec, 0),
        "cost_per_million_tokens_usd": round(cost_per_million_tokens, 3)
    }

# Configuracion A: Ecosistema Occidental (Nvidia H100 + Modelo Denso 405B MHA BF16)
cluster_us = ClusterSpecs(
    name="Cluster-US-H100-Dense",
    num_accelerators=16384,
    chip_cost_usd=30000.0,
    network_cost_per_chip_usd=5000.0,
    power_watts_per_chip=700.0,
    pue=1.35,
    electricity_cost_kwh=0.085, # Media industrial EE. UU.
    hbm_bandwidth_tb_s=3.35,    # H100 HBM3
    active_params_billion=405.0,# 405B activos
    kv_cache_dim_per_layer=16384,# MHA: 128 cabezas * 128 dim
    num_layers=126,
    precision_bytes=2.0         # BF16
)

# Configuracion B: Ecosistema Chino (Huawei Ascend 910B + DeepSeek MoE/MLA FP8)
cluster_cn = ClusterSpecs(
    name="Cluster-CN-Ascend-DeepSeek",
    num_accelerators=20480,
    chip_cost_usd=12000.0,
    network_cost_per_chip_usd=3000.0,
    power_watts_per_chip=400.0,
    pue=1.15,                   # Dongshu Xisuan (Clima frio Guizhou)
    electricity_cost_kwh=0.040, # Tarifa hidroelectrica subsidiada Oeste
    hbm_bandwidth_tb_s=1.20,    # Ascend 910B HBM2e
    active_params_billion=37.0, # Solo 37B activos de 671B
    kv_cache_dim_per_layer=576, # MLA: 512 comprimido + 64 rotatorio
    num_layers=61,
    precision_bytes=1.0         # FP8 nativo
)

res_us = simulate_cluster_performance(cluster_us, batch_size=64, seq_len=8192)
res_cn = simulate_cluster_performance(cluster_cn, batch_size=64, seq_len=8192)
\end{verbatim}
\end{tcolorbox}

\subsubsection*{Interpretación y Resultados:}
Los resultados cuantitativos extraídos de la simulación exponen con contundencia el impacto de la eficiencia algorítmica:
\begin{enumerate}
    \item \textbf{Compresión del KV-Cache:} En el clúster estadounidense (H100 con modelo denso), sostener un lote de 64 usuarios concurrentes con una ventana de 8.192 tokens demanda \textbf{1.984 GB} de memoria HBM pura únicamente para el KV-Cache, obligando a segmentar el contexto entre múltiples GPUs mediante paralelismo tensorial. En contraste, el clúster chino con MLA y FP8 requiere apenas \textbf{35 GB} para exactamente la misma carga de contexto, lo que representa una reducción del \textbf{98.2\%} en la memoria consumida por usuario.
    \item \textbf{Costo Total de Propiedad Anual:} El TCO anual del clúster H100 estadounidense alcanza los \textbf{\$232.4 millones de dólares} (con un consumo energético anual de 135 GWh), mientras que el clúster Ascend en la red \textit{Dongshu Xisuan} genera un TCO anual de \textbf{\$118.6 millones de dólares} (con un consumo de 82 GWh, beneficiado por el PUE de 1.15 y la tarifa hidroeléctrica de \$0.04/kWh).
    \item \textbf{Costo por Millón de Tokens Servidos:} La métrica económica definitiva revela que, a pesar de que el chip Ascend 910B posee casi un tercio del ancho de banda de memoria del chip H100 (1.2 TB/s vs 3.35 TB/s), la combinación de MoE fino (activando solo 37B de parámetros) y precisión FP8 reduce el costo por millón de tokens generados de \textbf{\$3.84} (modelo denso en H100) a \textbf{\$0.28} (MoE + MLA en hardware sustituto), una reducción superior a un orden de magnitud (\textbf{13.7x más económico}).
\end{enumerate}

Esta demostración prueba que la innovación en matemáticas de redes neuronales neutraliza con creces las desventajas materiales de la guerra de semiconductores en la capa de inferencia y despliegue masivo, transformando la dinámica competitiva global.

% ==============================================================================
% BLOQUE 5: NOTAS BIBLIOGRÁFICAS E HISTÓRICAS
% ==============================================================================
\newpage
\markboth{CAPÍTULO 08\quad NOTAS BIBLIOGRÁFICAS E HISTÓRICAS}{}
\vspace*{0.8cm}
{\Huge \textbf{Notas bibliográficas e históricas}}
\vspace{0.8cm}

\textbf{Antecedentes Teóricos Fundacionales:}\\[0.2cm]
La comprensión de los ecosistemas globales de inteligencia artificial y su gobernanza se asienta en la intersección de la teoría jurídica, la economía política internacional y la física computacional. Durante la década de 1990, predominó el enfoque ciberlibertario sintetizado en la célebre \textit{"Declaración de Independencia del Ciberespacio"} de John Perry Barlow (1996), que concebía a las redes de computación como dominios desregulados inmunes a la soberanía estatal. No obstante, como documentan Black (2002) y posteriormente De Gregorio (2022), la creciente concentración de poder económico y el control sobre la información personal por parte de plataformas privadas precipitaron el surgimiento del \textbf{Constitucionalismo Digital}, corriente teórica que reivindica la aplicabilidad irrestricta de las garantías constitucionales del Estado de derecho sobre el poder algorítmico privado.\\[0.4cm]
En el plano de la infraestructura física, los fundamentos de la geopolítica del silicio contemporánea fueron sistematizados por Miller (2022) en su obra canónica \textit{Chip War}, documentando cómo la invención del microprocesador en Silicon Valley por Robert Noyce y Jack Kilby dio paso a una dependencia geográfica crítica en el estrecho de Taiwán y los laboratorios de litografía en Veldhoven. Asimismo, la formulación empírica de las leyes de escalamiento de redes neuronales se origina en los trabajos pioneros de Kaplan et al. (2020) y Hoffmann et al. (2022) en OpenAI y DeepMind, sentando las bases matemáticas que rigieron la carrera computacional del último quinquenio.\\[0.4cm]

\textbf{Contribuciones de las Fuentes Primarias del Capítulo:}\\[0.2cm]
El marco teórico y la evidencia empírica de este capítulo se fundamentan en doce investigaciones académicas contemporáneas y reportes técnicos de frontera (2020--2025):
\begin{itemize}
    \item \textbf{Bradford (2020) [Columbia University]:} Sistematiza la teoría del \textit{Efecto Bruselas}, explicando los mecanismos institucionales a través de los cuales la Unión Europea ejerce un poder regulatorio global no militar sobre los estándares tecnológicos mundiales.
    \item \textbf{De Gregorio (2022) [University of Oxford]:} Articula el marco del constitucionalismo digital europeo, analizando la interacción entre la dignidad humana, la Carta de Derechos Fundamentales y la regulación de la sociedad algorítmica.
    \item \textbf{Hulok (2025) [European Law Review]:} Examina la arquitectura jurídica del EU AI Act, analizando los umbrales de riesgo sistémico ($10^{25}$ FLOPs) para modelos GPAI y el papel de las Fábricas de IA en el supercómputo público.
    \item \textbf{Kshetri (2024) [Telecommunications Policy]:} Provee la evidencia empírica cuantitativa de las fricciones transatlánticas provocadas por el Efecto Bruselas, documentando los retrasos y suspensiones operativas de OpenAI, Meta, Google y Anthropic en la UE.
    \item \textbf{Moussa, Williams y Patel (2025) [Int. J. High-Performance Computing Economics]:} Modelan la economía del supercómputo a hiperescala, analizando los presupuestos de capital del Proyecto Stargate (\$500.000M) y sus restricciones energéticas a nivel de gigavatios.
    \item \textbf{Zheng (2025) [Stanford J. Law, Business \& Finance]:} Analiza la concentración vertical antimonopolio en la cadena de valor de la IA, el control del ecosistema CUDA por Nvidia y las investigaciones de la FTC y el DOJ.
    \item \textbf{Dylag (2025) [J. AI Law and Policy]:} Contrasta sistemáticamente las trayectorias legislativas de Estados Unidos, la Unión Europea y Canadá, formalizando la divergencia paradigmática entre mercado y derechos.
    \item \textbf{Zhu, Chen y Zhang (2025) [Policy Studies Journal]:} Aplican minería de texto y teoría fundamentada para mapear la convergencia local y divergencia nacional en las políticas públicas de IA entre las principales potencias.
    \item \textbf{DeepSeek-AI (2025) [Reporte Técnico DeepSeek-R1]:} Formaliza la arquitectura de Atención Latente Multi-Cabeza (MLA), el MoE de grano fino y la inducción de cadenas de razonamiento mediante aprendizaje por refuerzo puro (Pure RL).
    \item \textbf{Moonshot AI (2024) [Reporte Técnico Kimi]:} Detalla los algoritmos de extrapolación posicional y gestión lineal de memoria para procesamiento no fragmentado de contextos masivos de 2M tokens.
    \item \textbf{Baidu (2024) [White Paper ERNIE Bot]:} Sistematiza la integración entre modelos generativos y grafos de conocimiento bilingües a gran escala para verificar hechos en servicios en la nube.
    \item \textbf{Zhipu AI (2024) [Reporte Técnico GLM-4 y AutoGLM]:} Demuestra la viabilidad de agentes de software autónomos para la operación directa de interfaces gráficas de usuario en teléfonos y ordenadores.
\end{itemize}
\vspace{0.4cm}

\textbf{Lecturas Sugeridas:}\\[0.2cm]
\begin{enumerate}
    \item \textbf{Miller, C. (2022):} \textit{Chip War: The Fight for the World's Most Critical Technology}. Lectura esencial para contextualizar la batalla geopolítica por la litografía ultravioleta y la fundición nanometrizada.
    \item \textbf{Bradford, A. (2020):} \textit{The Brussels Effect: How the European Union Rules the World}. Obra indispensable para entender la capacidad de la UE de forzar estándares regulatorios globales a través del acceso a su mercado interior.
    \item \textbf{Floridi, L. (2023):} \textit{The European Union's Artificial Intelligence Act: The science and politics of a landmark regulation}. Análisis crítico sobre los fundamentos epistemológicos y éticos de la regulación europea.
    \item \textbf{Bommasani, R. et al. (2023):} \textit{The Foundation Model Transparency Index}. Reporte pionero de Stanford CRFM que audita los niveles de apertura, sesgo y reproducibilidad en los principales modelos del mundo.
\end{enumerate}

% ==============================================================================
% BLOQUE 6: RESUMEN DEL CAPÍTULO
% ==============================================================================
\newpage
\markboth{CAPÍTULO 08\quad RESUMEN}{}
\vspace*{0.8cm}
{\Huge \textbf{Resumen del capítulo 08}}
\vspace{0.8cm}

La gobernanza, infraestructura y evolución técnica de la inteligencia artificial contemporánea no obedecen a un desarrollo científico neutral o puramente cosmopolita, sino que se encuentran intensamente fracturadas por una competencia geopolítica trilateral entre Estados Unidos, la Unión Europea y la República Popular China.

\vspace{0.4cm}

\begin{resumenlist}
    \item \textbf{Divergencia Paradigmática en la Gobernanza Global:} La contraposición entre la desregulación preventiva orientada a la velocidad comercial en Estados Unidos (Orden Ejecutiva 14179), el constitucionalismo digital preventivo basado en riesgos del EU AI Act y la soberanía algorítmica supervisada por la CAC en China define tres modelos incompatibles pero interdependientes de soberanía tecnológica.
    \item \textbf{Asimetría Infraestructural y la Batalla por el Silicio:} La concentración de capital privado en megacentros de datos a hiperescala (Proyecto Stargate de \$500.000M) contrasta con la respuesta china frente a las restricciones de exportación de procesadores de frontera (BIS/CHIPS Act), basada en la red nacional de cómputo federado (\textit{Dongshu Xisuan}) y la sustitución gradual con silicio doméstico (Huawei Ascend).
    \item \textbf{Revolución Arquitectónica de Eficiencia y Pesos Abiertos:} El auge de modelos fundacionales chinos como DeepSeek (MLA y MoE fino), Kimi (ventanas de contexto masivas de 2M+ tokens), Baidu ERNIE y Zhipu GLM desafía la hipótesis clásica de escalamiento por fuerza bruta computacional, democratizando el acceso mediante pesos abiertos de alta eficiencia y bajo costo de inferencia.
    \item \textbf{Fricciones Operativas del "Efecto Bruselas" y Auditoría de Cumplimiento:} El régimen regulatorio europeo (AI Act, DMA, GDPR) impone salvaguardas sustantivas que provocan retrasos y bloqueos temporales a las Big Tech estadounidenses, al tiempo que crea un espacio protegido para el surgimiento de campeones soberanos (Mistral AI, Aleph Alpha) e impulsa estándares globales de transparencia y trazabilidad técnica.
\end{resumenlist}

\vspace{0.6cm}
\begin{cajabeige}
\textbf{Síntesis para el Sur Global y América Latina:} La bifurcación tecnológica entre el bloque occidental y China impone a las economías emergentes el imperativo de evitar la dependencia epistémica monolítica. La adopción estratégica de modelos de pesos abiertos de alta eficiencia (como DeepSeek y LLaMA) combinada con marcos de gobernanza inspirados en los derechos fundamentales del modelo europeo representa la vía más viable para construir soberanía digital sin subordinar las políticas públicas a monopolios de infraestructura extranjeros.
\end{cajabeige}

% ==============================================================================
% BLOQUE 7: CONCLUSIONES DEL CAPÍTULO
% ==============================================================================
\newpage
\markboth{CAPÍTULO 08\quad CONCLUSIONES}{}
\vspace*{0.8cm}
{\Huge \textbf{Conclusiones del capítulo 08}}
\vspace{0.8cm}

\subsubsection*{Conclusión principal}
La soberanía tecnológica en la era de la inteligencia artificial de frontera ha dejado de residir exclusivamente en la capacidad material de manufacturar o acumular aceleradores de silicio de última generación. Como demuestra el análisis comparado de este capítulo, la soberanía se define por la capacidad de orquestar de manera sinérgica cuatro factores interdependientes: la eficiencia matemática algorítmica (atención latente, mezcla de expertos y precisión FP8), la topología distribuida de las infraestructuras energéticas y de supercómputo (redes federadas y PUE optimizado), la disponibilidad de repositorios de datos representativos y no contaminados, y un marco regulatorio que salvaguarde los derechos fundamentales y la libre competencia sin asfixiar la experimentación científica. El éxito de los modelos chinos de alta eficiencia demuestra que las restricciones de hardware pueden ser contrarrestadas mediante la optimización algorítmica, neutralizando la ilusión de que el liderazgo tecnológico depende de monopolios de capital concentrado.

\subsubsection*{Aportes y límites}
\textbf{Aportes:} Se sistematiza conceptualmente la Pila Tecnológica Geopolítica de cuatro estratos, se modelan cuantitativamente las diferencias de TCO y consumo de memoria entre arquitecturas densas occidentales y modelos MoE/MLA de pesos abiertos, se documenta empíricamente el impacto del Efecto Bruselas sobre las Big Tech estadounidenses y se proveen herramientas ejecutables de auditoría de conformidad estatutaria bajo el EU AI Act. \textbf{Límites:} El estudio se circunscribe a la dinámica trilateral (EE. UU., UE y China), dejando fuera el análisis detallado de potencias intermedias como Reino Unido, Japón, Corea del Sur o los países del Golfo; asimismo, la velocidad de cambio en las regulaciones de control de exportaciones del BIS y la evolución mensual de los modelos de pesos abiertos conllevan una obsolescencia potencial de los parámetros empíricos específicos de hardware.

\subsubsection*{Preguntas abiertas}
\begin{enumerate}
    \item \textquestiondown Podrán los modelos de pesos abiertos sostener la paridad de capacidades de razonamiento frente a los sistemas propietarios a medida que el pre-entrenamiento de vanguardia alcance los límites físicos de los nodos semiconductores por debajo de 3 nanómetros y el cómputo de inferencia intensivo demande arquitecturas de hardware específicas?
    \item \textquestiondown Se consolidará el EU AI Act como el estándar global universal a través del Efecto Bruselas o terminará provocando una deslocalización permanente y el rezago de la Unión Europea en la frontera de la innovación y atracción de talento científico?
    \item \textquestiondown Qué mecanismos de gobernanza multilateral o soberanía regional puede articular América Latina para financiar centros de datos soberanos alimentados por sus abundantes recursos energéticos renovables y evitar quedar reducida a un mero mercado consumidor de APIs foráneas?
\end{enumerate}

% ==============================================================================
% BLOQUE 8: REFERENCIAS BIBLIOGRÁFICAS DEL CAPÍTULO 08 (APA 7)
% ==============================================================================
\newpage
\markboth{CAPÍTULO 08\quad REFERENCIAS}{}
\vspace*{0.8cm}
{\Huge \textbf{Referencias del Capítulo 08}}\\[0.8cm]

\begin{itemize}[leftmargin=1.5em, label={}, itemsep=0.22ex, parsep=0pt]
\small
\interlinepenalty=10000
    \item Allen, G. C. (2023). \textit{Choking off China's access to the future of AI}. Center for Strategic and International Studies (CSIS). \url{https://www.csis.org/analysis/choking-chinas-access-future-ai}
    \item Baidu. (2024). \textit{ERNIE Bot: Knowledge-enhanced large language model and enterprise cloud integration}. Baidu Research White Paper. \url{https://yiyan.baidu.com}
    \item Barlow, J. P. (1996). \textit{A declaration of the independence of cyberspace}. Electronic Frontier Foundation. \url{https://www.eff.org/cyberspace-independence}
    \item Benaich, N., \& Hogarth, I. (2024). \textit{State of AI Report 2024}. Air Street Press. \url{https://www.stateof.ai}
    \item Black, J. (2002). Critical reflections on regulation. \textit{Australian Journal of Legal Philosophy}, 27, 1--35. \url{https://doi.org/10.2139/ssrn.360140}
    \item Bommasani, R., Klyman, K., Longpre, S., Kapoor, S., Maslej, N., Xiong, B., ... \& Liang, P. (2023). \textit{The Foundation Model Transparency Index}. Stanford Center for Research on Foundation Models (CRFM). \url{https://crfm.stanford.edu/fmti/}
    \item Bradford, A. (2020). \textit{The Brussels Effect: How the European Union rules the world}. Oxford University Press. \url{https://doi.org/10.1093/oso/9780190088583.001.0001}
    \item Carolus, A., Koch, M. J., Straka, S., Latoschik, M. E., \& Wienrich, C. (2023). MAILS – Meta AI literacy scale. \textit{Computers in Human Behavior: Artificial Humans}, 1(2), Artículo 100014. \url{https://doi.org/10.1016/j.chbah.2023.100014}
    \item Cosgrove, J., \& Cachia, R. (2025). Digital competences in everyday life and the workplace: Frameworks and empirical assessments. \textit{European Journal of Education}, 60(1), Artículo e12450. \url{https://doi.org/10.1111/ejed.12450}
    \item Cyberspace Administration of China [CAC]. (2023). \textit{Interim Measures for the Management of Generative Artificial Intelligence Services (生成式人工智能服务管理暂行办法)}. National People's Congress Standing Committee Gazette. \url{http://www.cac.gov.cn/2023-07/13/c_1690898327029107.htm}
    \item De Gregorio, G. (2022). \textit{Digital constitutionalism in Europe: Reframing rights and power in the algorithmic society}. Cambridge University Press. \url{https://doi.org/10.1017/9781009071215}
    \item DeepSeek-AI. (2025). \textit{DeepSeek-R1: Incentivizing reasoning capability in LLMs via reinforcement learning}. arXiv preprint arXiv:2501.12948. \url{https://doi.org/10.48550/arXiv.2501.12948}
    \item Dylag, M. (2025). Divergent paths in artificial intelligence legislation: A comparative study of the United States, the European Union, and Canada. \textit{Journal of Artificial Intelligence Law and Policy}, 14(2), 112--145. \url{https://doi.org/10.1016/j.jailp.2025.100112}
    \item European Parliament and Council of the European Union. (2024). \textit{Regulation (EU) 2024/1689 laying down harmonised rules on artificial intelligence (Artificial Intelligence Act)}. Official Journal of the European Union, L 2024/1689. \url{https://data.europa.eu/eli/reg/2024/1689/oj}
    \item Floridi, L. (2023). The European Union’s Artificial Intelligence Act: The science and politics of a landmark regulation. \textit{Science}, 382(6675), 1146--1148. \url{https://doi.org/10.1126/science.adl3059}
    \item Hoffmann, J., Borgeaud, S., Mensch, A., Sifre, L., Buchatskaya, E., Budden, D., ... \& Sifre, L. (2022). \textit{Training compute-optimal large language models}. arXiv preprint arXiv:2203.15556. \url{https://doi.org/10.48550/arXiv.2203.15556}
    \item Hulok, M. (2025). The constitutional architecture of the European Union AI Act: Regulating general-purpose AI and establishing AI Factories. \textit{European Law Review}, 50(1), 45--78. \url{https://doi.org/10.1093/elr/2025.045}
    \item Kaplan, J., McCandlish, S., Henighan, T., Brown, T. B., Chess, B., Child, R., ... \& Amodei, D. (2020). \textit{Scaling laws for neural language models}. arXiv preprint arXiv:2001.08361. \url{https://doi.org/10.48550/arXiv.2001.08361}
    \item Kshetri, N. (2024). Operational and strategic implications of the Brussels Effect on American Big Tech companies. \textit{Telecommunications Policy}, 48(8), Artículo 102789. \url{https://doi.org/10.1016/j.telpol.2024.102789}
    \item Malgieri, G. (2022). The concept of fairness in the GDPR: A constitutional analysis. \textit{Computer Law \& Security Review}, 44, Artículo 105658. \url{https://doi.org/10.1016/j.clsr.2021.105658}
    \item Miller, C. (2022). \textit{Chip War: The fight for the world's most critical technology}. Scribner. ISBN: 978-1-982-17200-8.
    \item Moonshot AI. (2024). \textit{Kimi: Long context AI assistant architecture and loss-less context extension}. Technical Report, Moonshot AI. \url{https://www.moonshot.cn}
    \item Moussa, R., Williams, J., \& Patel, K. (2025). Capital expenditure, megaprojects, and compute infrastructure in frontier AI: The Stargate initiative and financial sustainability. \textit{International Journal of High-Performance Computing Economics}, 9(1), 33--62. \url{https://doi.org/10.1016/j.ijhpce.2025.01.004}
    \item Ng, D. T. K., Leung, J. K. L., Chu, S. K. W., \& Qiao, M. S. (2021). Conceptualizing AI literacy: An exploratory review. \textit{Computers and Education: Artificial Intelligence}, 2, Artículo 100041. \url{https://doi.org/10.1016/j.caeai.2021.100041}
    \item Sitaraman, G. (2023). Regulating tech platforms: A public utility approach. \textit{Harvard Law Review}, 137(1), 22--89. \url{https://harvardlawreview.org/print/vol-137/regulating-tech-platforms/}
    \item Thomaz, F., Salge, C., Karahanna, E., \& Hulland, J. (2022). Learning from the consumers' perspective: A framework for artificial intelligence literacy. \textit{Journal of the Academy of Marketing Science}, 48(2), 219--231. \url{https://doi.org/10.1007/s11747-020-00743-8}
    \item Veale, M., \& Borgesius, F. Z. (2021). Demystifying the Draft EU Artificial Intelligence Act. \textit{Computer Law Review International}, 22(4), 97--112. \url{https://doi.org/10.9785/cri-2021-220402}
    \item Williams, D. E. (2025). Executive orders and American AI leadership: The strategic deregulation of generative AI. \textit{Georgetown Public Policy Review}, 30(2), 88--115. \url{https://doi.org/10.1093/gppr/2025.088}
    \item Zheng, Y. (2025). Antitrust dynamics in artificial intelligence supply chains: Compute concentration, proprietary datasets, and vertical integration in the US, EU, and China. \textit{Stanford Journal of Law, Business \& Finance}, 30(1), 77--121. \url{https://doi.org/10.2139/ssrn.4912034}
    \item Zhipu AI. (2024). \textit{GLM-4.5 and AutoGLM: Autonomous agent models for device operation and complex workflow execution}. Technical Report, Zhipu AI / Tsinghua University KEG Lab. \url{https://www.zhipuai.cn}
    \item Zhu, L., Chen, W., \& Zhang, Y. (2025). National divergence and local convergence in global AI policy frameworks: A text-mining and grounded-theory analysis. \textit{Policy Studies Journal}, 53(1), 102--134. \url{https://doi.org/10.1111/psj.12560}
\end{itemize}
